% paper30-semantic-control.tex
% Semantic Control Laws: Feedback-Driven Verification Orchestration
% Paper 30 of the Judgment Geometry series.
% Compile: pdflatex paper30-semantic-control

\documentclass[11pt]{article}
% jugeo-common.tex — shared preamble for all Judgment Geometry papers
% Include via: % jugeo-common.tex — shared preamble for all Judgment Geometry papers
% Include via: % jugeo-common.tex — shared preamble for all Judgment Geometry papers
% Include via: \input{jugeo-common}

% ── Packages ────────────────────────────────────────────────────────────
\usepackage{amsmath,amssymb,amsthm,mathtools}
\usepackage{tikz-cd}
\usepackage{listings}
\usepackage{xcolor}
\usepackage{xspace}
\usepackage{booktabs}
\usepackage{multirow}
\usepackage{enumitem}
\usepackage[expansion=false]{microtype}
\usepackage{hyperref}
\usepackage{cleveref}
% \usepackage{thmtools}  % disabled: not available in this TeX installation
\usepackage{float}
\usepackage{caption}
\usepackage{subcaption}
\usepackage{tabularx}
\usepackage{stmaryrd}       % \llbracket, \rrbracket
\usepackage{bbm}            % \mathbbm{1}

% ── Page geometry ───────────────────────────────────────────────────────
\usepackage[margin=1in]{geometry}

% ── Theorem environments ────────────────────────────────────────────────
\theoremstyle{plain}
\newtheorem{theorem}{Theorem}[section]
\newtheorem{lemma}[theorem]{Lemma}
\newtheorem{proposition}[theorem]{Proposition}
\newtheorem{corollary}[theorem]{Corollary}

\theoremstyle{definition}
\newtheorem{definition}[theorem]{Definition}
\newtheorem{example}[theorem]{Example}
\newtheorem{construction}[theorem]{Construction}

\theoremstyle{remark}
\newtheorem{remark}[theorem]{Remark}
\newtheorem{notation}[theorem]{Notation}

% ── Python listings style ──────────────────────────────────────────────
\definecolor{jugeoBlue}{HTML}{2563EB}
\definecolor{jugeoGreen}{HTML}{059669}
\definecolor{jugeoRed}{HTML}{DC2626}
\definecolor{jugeoPurple}{HTML}{7C3AED}
\definecolor{jugeoGray}{HTML}{6B7280}
\definecolor{jugeoBg}{HTML}{F8FAFC}

\lstdefinestyle{jugeo-python}{
  language=Python,
  basicstyle=\ttfamily\small,
  keywordstyle=\color{jugeoBlue}\bfseries,
  stringstyle=\color{jugeoGreen},
  commentstyle=\color{jugeoGray}\itshape,
  numberstyle=\tiny\color{jugeoGray},
  backgroundcolor=\color{jugeoBg},
  numbers=left,
  numbersep=6pt,
  frame=single,
  framerule=0.4pt,
  rulecolor=\color{jugeoGray!40},
  breaklines=true,
  breakatwhitespace=true,
  showstringspaces=false,
  tabsize=4,
  xleftmargin=1.5em,
  framexleftmargin=1.5em,
  morekeywords={dataclass,frozen,field,Enum,IntEnum,auto,Any,
                Mapping,Sequence,Optional,match,case},
  literate={→}{{$\to$}}1 {←}{{$\leftarrow$}}1
           {≤}{{$\leq$}}1 {≥}{{$\geq$}}1 {≠}{{$\neq$}}1
           {∈}{{$\in$}}1 {∀}{{$\forall$}}1 {∃}{{$\exists$}}1
           {⊕}{{$\oplus$}}1 {⊖}{{$\ominus$}}1,
  escapeinside={(*@}{@*)},
}
\lstset{style=jugeo-python}

\lstdefinestyle{jugeo-output}{
  basicstyle=\ttfamily\small,
  backgroundcolor=\color{jugeoBg},
  frame=single,
  framerule=0.4pt,
  rulecolor=\color{jugeoGray!40},
  breaklines=true,
  showstringspaces=false,
  xleftmargin=1.5em,
  framexleftmargin=0.5em,
  numbers=none,
}

% ── JuGeo-specific macros ──────────────────────────────────────────────

% Judgment 8-tuple
\newcommand{\J}{\mathcal{J}}
\newcommand{\Jud}[8]{(#1,\,#2,\,#3,\,#4,\,#5,\,#6,\,#7,\,#8)}
\newcommand{\JudShort}[1]{\J_{#1}}

% Components
\newcommand{\coord}{\mathsf{c}}            % coordinate
\newcommand{\prop}{\varphi}                 % proposition
\newcommand{\carrier}{\mathcal{A}}          % carrier type
\newcommand{\evid}{\mathcal{E}}             % evidence bundle
\newcommand{\oblig}{\mathcal{O}}            % obligations
\newcommand{\obstr}{\mathcal{B}}            % obstructions
\newcommand{\trust}{\mathcal{T}}            % trust annotation
\newcommand{\prov}{\Pi}                     % provenance

% Trust algebra
\newcommand{\Talg}{\mathfrak{T}}            % trust ordered algebra
\newcommand{\Eadm}{\mathcal{E}_{\mathrm{adm}}}
\newcommand{\tleq}{\preceq}                 % trust ordering
\newcommand{\tjoin}{\oplus}                 % conservative join
\newcommand{\tatten}{\ominus}               % attenuation
\newcommand{\tpromote}{\uparrow_\pi}        % promotion
\newcommand{\tdemote}{\downarrow_\chi}       % demotion

% Trust tiers
\newcommand{\Tcontra}{\ensuremath{\mathsf{CONTRADICTED}}}
\newcommand{\Tunver}{\ensuremath{\mathsf{UNVERIFIED}}}
\newcommand{\Tcopilot}{\ensuremath{\mathsf{COPILOT}}}
\newcommand{\Toracle}{\ensuremath{\mathsf{ORACLE}}}
\newcommand{\Truntime}{\ensuremath{\mathsf{RUNTIME}}}
\newcommand{\Tsolver}{\ensuremath{\mathsf{SOLVER}}}
\newcommand{\Tproof}{\ensuremath{\mathsf{PROOF}}}

% Site and geometry
\newcommand{\Site}{\mathcal{S}}             % semantic site
\newcommand{\Cov}{\mathrm{Cov}}            % covering family
\newcommand{\Ob}{\mathrm{Ob}}              % objects
\newcommand{\Mor}{\mathrm{Mor}}            % morphisms
\newcommand{\res}{\mathrm{res}}            % restriction
\newcommand{\Cech}{\check{C}}              % Čech complex
\newcommand{\Hcoh}[1]{H^{#1}}             % cohomology group

% Descent
\newcommand{\descent}{\mathrm{desc}}
\newcommand{\glue}{\mathrm{glue}}
\newcommand{\overlap}[2]{#1 \cap #2}

% Semantic moves
\newcommand{\VERIFY}{\textsc{Verify}}
\newcommand{\CONSTRUCT}{\textsc{Construct}}
\newcommand{\NORMALIZE}{\textsc{Normalize}}
\newcommand{\GLUE}{\textsc{Glue}}
\newcommand{\RESTRICT}{\textsc{Restrict}}
\newcommand{\TRANSPORT}{\textsc{Transport}}
\newcommand{\REPAIR}{\textsc{Repair}}
\newcommand{\PROMOTE}{\textsc{Promote}}
\newcommand{\CHALLENGE}{\textsc{Challenge}}
\newcommand{\ESCALATE}{\textsc{Escalate}}

% Fragments
\newcommand{\QFLIA}{\texttt{QF\_LIA}}
\newcommand{\QFLRA}{\texttt{QF\_LRA}}
\newcommand{\QFBV}{\texttt{QF\_BV}}
\newcommand{\QFUF}{\texttt{QF\_UF}}

% Misc
\newcommand{\Python}{\textsc{Python}}
\newcommand{\JuGeo}{\textsc{JuGeo}}
\newcommand{\LEAN}{\textsc{Lean}\,4}
\newcommand{\Fstar}{F$^{\star}$}
\newcommand{\Coq}{\textsc{Coq}}
\newcommand{\Dafny}{\textsc{Dafny}}

% Common shortcuts
\newcommand{\ie}{i.e.\@\xspace}
\newcommand{\eg}{e.g.\@\xspace}
\newcommand{\cf}{cf.\@\xspace}
\newcommand{\wrt}{w.r.t.\@\xspace}

% Evaluation shortcuts
\newcommand{\numcases}{300}
\newcommand{\numspec}{100}
\newcommand{\numeq}{100}
\newcommand{\numbug}{100}
\newcommand{\medtime}{6.5\,ms}
\newcommand{\totaltime}{1.949\,s}


% ── Packages ────────────────────────────────────────────────────────────
\usepackage{amsmath,amssymb,amsthm,mathtools}
\usepackage{tikz-cd}
\usepackage{listings}
\usepackage{xcolor}
\usepackage{xspace}
\usepackage{booktabs}
\usepackage{multirow}
\usepackage{enumitem}
\usepackage[expansion=false]{microtype}
\usepackage{hyperref}
\usepackage{cleveref}
% \usepackage{thmtools}  % disabled: not available in this TeX installation
\usepackage{float}
\usepackage{caption}
\usepackage{subcaption}
\usepackage{tabularx}
\usepackage{stmaryrd}       % \llbracket, \rrbracket
\usepackage{bbm}            % \mathbbm{1}

% ── Page geometry ───────────────────────────────────────────────────────
\usepackage[margin=1in]{geometry}

% ── Theorem environments ────────────────────────────────────────────────
\theoremstyle{plain}
\newtheorem{theorem}{Theorem}[section]
\newtheorem{lemma}[theorem]{Lemma}
\newtheorem{proposition}[theorem]{Proposition}
\newtheorem{corollary}[theorem]{Corollary}

\theoremstyle{definition}
\newtheorem{definition}[theorem]{Definition}
\newtheorem{example}[theorem]{Example}
\newtheorem{construction}[theorem]{Construction}

\theoremstyle{remark}
\newtheorem{remark}[theorem]{Remark}
\newtheorem{notation}[theorem]{Notation}

% ── Python listings style ──────────────────────────────────────────────
\definecolor{jugeoBlue}{HTML}{2563EB}
\definecolor{jugeoGreen}{HTML}{059669}
\definecolor{jugeoRed}{HTML}{DC2626}
\definecolor{jugeoPurple}{HTML}{7C3AED}
\definecolor{jugeoGray}{HTML}{6B7280}
\definecolor{jugeoBg}{HTML}{F8FAFC}

\lstdefinestyle{jugeo-python}{
  language=Python,
  basicstyle=\ttfamily\small,
  keywordstyle=\color{jugeoBlue}\bfseries,
  stringstyle=\color{jugeoGreen},
  commentstyle=\color{jugeoGray}\itshape,
  numberstyle=\tiny\color{jugeoGray},
  backgroundcolor=\color{jugeoBg},
  numbers=left,
  numbersep=6pt,
  frame=single,
  framerule=0.4pt,
  rulecolor=\color{jugeoGray!40},
  breaklines=true,
  breakatwhitespace=true,
  showstringspaces=false,
  tabsize=4,
  xleftmargin=1.5em,
  framexleftmargin=1.5em,
  morekeywords={dataclass,frozen,field,Enum,IntEnum,auto,Any,
                Mapping,Sequence,Optional,match,case},
  literate={→}{{$\to$}}1 {←}{{$\leftarrow$}}1
           {≤}{{$\leq$}}1 {≥}{{$\geq$}}1 {≠}{{$\neq$}}1
           {∈}{{$\in$}}1 {∀}{{$\forall$}}1 {∃}{{$\exists$}}1
           {⊕}{{$\oplus$}}1 {⊖}{{$\ominus$}}1,
  escapeinside={(*@}{@*)},
}
\lstset{style=jugeo-python}

\lstdefinestyle{jugeo-output}{
  basicstyle=\ttfamily\small,
  backgroundcolor=\color{jugeoBg},
  frame=single,
  framerule=0.4pt,
  rulecolor=\color{jugeoGray!40},
  breaklines=true,
  showstringspaces=false,
  xleftmargin=1.5em,
  framexleftmargin=0.5em,
  numbers=none,
}

% ── JuGeo-specific macros ──────────────────────────────────────────────

% Judgment 8-tuple
\newcommand{\J}{\mathcal{J}}
\newcommand{\Jud}[8]{(#1,\,#2,\,#3,\,#4,\,#5,\,#6,\,#7,\,#8)}
\newcommand{\JudShort}[1]{\J_{#1}}

% Components
\newcommand{\coord}{\mathsf{c}}            % coordinate
\newcommand{\prop}{\varphi}                 % proposition
\newcommand{\carrier}{\mathcal{A}}          % carrier type
\newcommand{\evid}{\mathcal{E}}             % evidence bundle
\newcommand{\oblig}{\mathcal{O}}            % obligations
\newcommand{\obstr}{\mathcal{B}}            % obstructions
\newcommand{\trust}{\mathcal{T}}            % trust annotation
\newcommand{\prov}{\Pi}                     % provenance

% Trust algebra
\newcommand{\Talg}{\mathfrak{T}}            % trust ordered algebra
\newcommand{\Eadm}{\mathcal{E}_{\mathrm{adm}}}
\newcommand{\tleq}{\preceq}                 % trust ordering
\newcommand{\tjoin}{\oplus}                 % conservative join
\newcommand{\tatten}{\ominus}               % attenuation
\newcommand{\tpromote}{\uparrow_\pi}        % promotion
\newcommand{\tdemote}{\downarrow_\chi}       % demotion

% Trust tiers
\newcommand{\Tcontra}{\ensuremath{\mathsf{CONTRADICTED}}}
\newcommand{\Tunver}{\ensuremath{\mathsf{UNVERIFIED}}}
\newcommand{\Tcopilot}{\ensuremath{\mathsf{COPILOT}}}
\newcommand{\Toracle}{\ensuremath{\mathsf{ORACLE}}}
\newcommand{\Truntime}{\ensuremath{\mathsf{RUNTIME}}}
\newcommand{\Tsolver}{\ensuremath{\mathsf{SOLVER}}}
\newcommand{\Tproof}{\ensuremath{\mathsf{PROOF}}}

% Site and geometry
\newcommand{\Site}{\mathcal{S}}             % semantic site
\newcommand{\Cov}{\mathrm{Cov}}            % covering family
\newcommand{\Ob}{\mathrm{Ob}}              % objects
\newcommand{\Mor}{\mathrm{Mor}}            % morphisms
\newcommand{\res}{\mathrm{res}}            % restriction
\newcommand{\Cech}{\check{C}}              % Čech complex
\newcommand{\Hcoh}[1]{H^{#1}}             % cohomology group

% Descent
\newcommand{\descent}{\mathrm{desc}}
\newcommand{\glue}{\mathrm{glue}}
\newcommand{\overlap}[2]{#1 \cap #2}

% Semantic moves
\newcommand{\VERIFY}{\textsc{Verify}}
\newcommand{\CONSTRUCT}{\textsc{Construct}}
\newcommand{\NORMALIZE}{\textsc{Normalize}}
\newcommand{\GLUE}{\textsc{Glue}}
\newcommand{\RESTRICT}{\textsc{Restrict}}
\newcommand{\TRANSPORT}{\textsc{Transport}}
\newcommand{\REPAIR}{\textsc{Repair}}
\newcommand{\PROMOTE}{\textsc{Promote}}
\newcommand{\CHALLENGE}{\textsc{Challenge}}
\newcommand{\ESCALATE}{\textsc{Escalate}}

% Fragments
\newcommand{\QFLIA}{\texttt{QF\_LIA}}
\newcommand{\QFLRA}{\texttt{QF\_LRA}}
\newcommand{\QFBV}{\texttt{QF\_BV}}
\newcommand{\QFUF}{\texttt{QF\_UF}}

% Misc
\newcommand{\Python}{\textsc{Python}}
\newcommand{\JuGeo}{\textsc{JuGeo}}
\newcommand{\LEAN}{\textsc{Lean}\,4}
\newcommand{\Fstar}{F$^{\star}$}
\newcommand{\Coq}{\textsc{Coq}}
\newcommand{\Dafny}{\textsc{Dafny}}

% Common shortcuts
\newcommand{\ie}{i.e.\@\xspace}
\newcommand{\eg}{e.g.\@\xspace}
\newcommand{\cf}{cf.\@\xspace}
\newcommand{\wrt}{w.r.t.\@\xspace}

% Evaluation shortcuts
\newcommand{\numcases}{300}
\newcommand{\numspec}{100}
\newcommand{\numeq}{100}
\newcommand{\numbug}{100}
\newcommand{\medtime}{6.5\,ms}
\newcommand{\totaltime}{1.949\,s}


% ── Packages ────────────────────────────────────────────────────────────
\usepackage{amsmath,amssymb,amsthm,mathtools}
\usepackage{tikz-cd}
\usepackage{listings}
\usepackage{xcolor}
\usepackage{xspace}
\usepackage{booktabs}
\usepackage{multirow}
\usepackage{enumitem}
\usepackage[expansion=false]{microtype}
\usepackage{hyperref}
\usepackage{cleveref}
% \usepackage{thmtools}  % disabled: not available in this TeX installation
\usepackage{float}
\usepackage{caption}
\usepackage{subcaption}
\usepackage{tabularx}
\usepackage{stmaryrd}       % \llbracket, \rrbracket
\usepackage{bbm}            % \mathbbm{1}

% ── Page geometry ───────────────────────────────────────────────────────
\usepackage[margin=1in]{geometry}

% ── Theorem environments ────────────────────────────────────────────────
\theoremstyle{plain}
\newtheorem{theorem}{Theorem}[section]
\newtheorem{lemma}[theorem]{Lemma}
\newtheorem{proposition}[theorem]{Proposition}
\newtheorem{corollary}[theorem]{Corollary}

\theoremstyle{definition}
\newtheorem{definition}[theorem]{Definition}
\newtheorem{example}[theorem]{Example}
\newtheorem{construction}[theorem]{Construction}

\theoremstyle{remark}
\newtheorem{remark}[theorem]{Remark}
\newtheorem{notation}[theorem]{Notation}

% ── Python listings style ──────────────────────────────────────────────
\definecolor{jugeoBlue}{HTML}{2563EB}
\definecolor{jugeoGreen}{HTML}{059669}
\definecolor{jugeoRed}{HTML}{DC2626}
\definecolor{jugeoPurple}{HTML}{7C3AED}
\definecolor{jugeoGray}{HTML}{6B7280}
\definecolor{jugeoBg}{HTML}{F8FAFC}

\lstdefinestyle{jugeo-python}{
  language=Python,
  basicstyle=\ttfamily\small,
  keywordstyle=\color{jugeoBlue}\bfseries,
  stringstyle=\color{jugeoGreen},
  commentstyle=\color{jugeoGray}\itshape,
  numberstyle=\tiny\color{jugeoGray},
  backgroundcolor=\color{jugeoBg},
  numbers=left,
  numbersep=6pt,
  frame=single,
  framerule=0.4pt,
  rulecolor=\color{jugeoGray!40},
  breaklines=true,
  breakatwhitespace=true,
  showstringspaces=false,
  tabsize=4,
  xleftmargin=1.5em,
  framexleftmargin=1.5em,
  morekeywords={dataclass,frozen,field,Enum,IntEnum,auto,Any,
                Mapping,Sequence,Optional,match,case},
  literate={→}{{$\to$}}1 {←}{{$\leftarrow$}}1
           {≤}{{$\leq$}}1 {≥}{{$\geq$}}1 {≠}{{$\neq$}}1
           {∈}{{$\in$}}1 {∀}{{$\forall$}}1 {∃}{{$\exists$}}1
           {⊕}{{$\oplus$}}1 {⊖}{{$\ominus$}}1,
  escapeinside={(*@}{@*)},
}
\lstset{style=jugeo-python}

\lstdefinestyle{jugeo-output}{
  basicstyle=\ttfamily\small,
  backgroundcolor=\color{jugeoBg},
  frame=single,
  framerule=0.4pt,
  rulecolor=\color{jugeoGray!40},
  breaklines=true,
  showstringspaces=false,
  xleftmargin=1.5em,
  framexleftmargin=0.5em,
  numbers=none,
}

% ── JuGeo-specific macros ──────────────────────────────────────────────

% Judgment 8-tuple
\newcommand{\J}{\mathcal{J}}
\newcommand{\Jud}[8]{(#1,\,#2,\,#3,\,#4,\,#5,\,#6,\,#7,\,#8)}
\newcommand{\JudShort}[1]{\J_{#1}}

% Components
\newcommand{\coord}{\mathsf{c}}            % coordinate
\newcommand{\prop}{\varphi}                 % proposition
\newcommand{\carrier}{\mathcal{A}}          % carrier type
\newcommand{\evid}{\mathcal{E}}             % evidence bundle
\newcommand{\oblig}{\mathcal{O}}            % obligations
\newcommand{\obstr}{\mathcal{B}}            % obstructions
\newcommand{\trust}{\mathcal{T}}            % trust annotation
\newcommand{\prov}{\Pi}                     % provenance

% Trust algebra
\newcommand{\Talg}{\mathfrak{T}}            % trust ordered algebra
\newcommand{\Eadm}{\mathcal{E}_{\mathrm{adm}}}
\newcommand{\tleq}{\preceq}                 % trust ordering
\newcommand{\tjoin}{\oplus}                 % conservative join
\newcommand{\tatten}{\ominus}               % attenuation
\newcommand{\tpromote}{\uparrow_\pi}        % promotion
\newcommand{\tdemote}{\downarrow_\chi}       % demotion

% Trust tiers
\newcommand{\Tcontra}{\ensuremath{\mathsf{CONTRADICTED}}}
\newcommand{\Tunver}{\ensuremath{\mathsf{UNVERIFIED}}}
\newcommand{\Tcopilot}{\ensuremath{\mathsf{COPILOT}}}
\newcommand{\Toracle}{\ensuremath{\mathsf{ORACLE}}}
\newcommand{\Truntime}{\ensuremath{\mathsf{RUNTIME}}}
\newcommand{\Tsolver}{\ensuremath{\mathsf{SOLVER}}}
\newcommand{\Tproof}{\ensuremath{\mathsf{PROOF}}}

% Site and geometry
\newcommand{\Site}{\mathcal{S}}             % semantic site
\newcommand{\Cov}{\mathrm{Cov}}            % covering family
\newcommand{\Ob}{\mathrm{Ob}}              % objects
\newcommand{\Mor}{\mathrm{Mor}}            % morphisms
\newcommand{\res}{\mathrm{res}}            % restriction
\newcommand{\Cech}{\check{C}}              % Čech complex
\newcommand{\Hcoh}[1]{H^{#1}}             % cohomology group

% Descent
\newcommand{\descent}{\mathrm{desc}}
\newcommand{\glue}{\mathrm{glue}}
\newcommand{\overlap}[2]{#1 \cap #2}

% Semantic moves
\newcommand{\VERIFY}{\textsc{Verify}}
\newcommand{\CONSTRUCT}{\textsc{Construct}}
\newcommand{\NORMALIZE}{\textsc{Normalize}}
\newcommand{\GLUE}{\textsc{Glue}}
\newcommand{\RESTRICT}{\textsc{Restrict}}
\newcommand{\TRANSPORT}{\textsc{Transport}}
\newcommand{\REPAIR}{\textsc{Repair}}
\newcommand{\PROMOTE}{\textsc{Promote}}
\newcommand{\CHALLENGE}{\textsc{Challenge}}
\newcommand{\ESCALATE}{\textsc{Escalate}}

% Fragments
\newcommand{\QFLIA}{\texttt{QF\_LIA}}
\newcommand{\QFLRA}{\texttt{QF\_LRA}}
\newcommand{\QFBV}{\texttt{QF\_BV}}
\newcommand{\QFUF}{\texttt{QF\_UF}}

% Misc
\newcommand{\Python}{\textsc{Python}}
\newcommand{\JuGeo}{\textsc{JuGeo}}
\newcommand{\LEAN}{\textsc{Lean}\,4}
\newcommand{\Fstar}{F$^{\star}$}
\newcommand{\Coq}{\textsc{Coq}}
\newcommand{\Dafny}{\textsc{Dafny}}

% Common shortcuts
\newcommand{\ie}{i.e.\@\xspace}
\newcommand{\eg}{e.g.\@\xspace}
\newcommand{\cf}{cf.\@\xspace}
\newcommand{\wrt}{w.r.t.\@\xspace}

% Evaluation shortcuts
\newcommand{\numcases}{300}
\newcommand{\numspec}{100}
\newcommand{\numeq}{100}
\newcommand{\numbug}{100}
\newcommand{\medtime}{6.5\,ms}
\newcommand{\totaltime}{1.949\,s}

% experiment-data.tex — AUTO-GENERATED from real jugeo CLI runs
% DO NOT EDIT — regenerate with: python3 experiments/run_universal_metrics.py
% Generated from 30 programs across 6 domains

\newcommand{\expTotalPrograms}{30}
\newcommand{\expVerified}{30}
\newcommand{\expAccuracy}{100.0\%}
\newcommand{\expCoordsMin}{2}
\newcommand{\expCoordsMax}{10}
\newcommand{\expCoordsMean}{5.2}
\newcommand{\expCoordsMedian}{5.5}
\newcommand{\expPropsMin}{4}
\newcommand{\expPropsMax}{20}
\newcommand{\expPropsMean}{10.4}
\newcommand{\expPropsSum}{311}
\newcommand{\expPropsOkSum}{310}
\newcommand{\expTimeMin}{3.506\,s}
\newcommand{\expTimeMax}{6.714\,s}
\newcommand{\expTimeMean}{4.64\,s}
\newcommand{\expTimeMedian}{4.508\,s}
\newcommand{\expTimeTotal}{139.204\,s}
\newcommand{\expObstructionTotal}{0}
\newcommand{\expHOneZero}{30}
\newcommand{\expDomainCount}{6}
\newcommand{\expTrustCOPILOTSUGGESTED}{30}
\newcommand{\expDomainDsCount}{5}
\newcommand{\expDomainDsVerified}{5}
\newcommand{\expDomainDsAvgCoords}{7.6}
\newcommand{\expDomainDsAvgProps}{15.2}
\newcommand{\expDomainMathCount}{5}
\newcommand{\expDomainMathVerified}{5}
\newcommand{\expDomainMathAvgCoords}{4.8}
\newcommand{\expDomainMathAvgProps}{9.4}
\newcommand{\expDomainSmCount}{5}
\newcommand{\expDomainSmVerified}{5}
\newcommand{\expDomainSmAvgCoords}{7.6}
\newcommand{\expDomainSmAvgProps}{15.2}
\newcommand{\expDomainSortCount}{5}
\newcommand{\expDomainSortVerified}{5}
\newcommand{\expDomainSortAvgCoords}{2.4}
\newcommand{\expDomainSortAvgProps}{4.8}
\newcommand{\expDomainStrCount}{5}
\newcommand{\expDomainStrVerified}{5}
\newcommand{\expDomainStrAvgCoords}{3.2}
\newcommand{\expDomainStrAvgProps}{6.4}
\newcommand{\expDomainWebCount}{5}
\newcommand{\expDomainWebVerified}{5}
\newcommand{\expDomainWebAvgCoords}{5.6}
\newcommand{\expDomainWebAvgProps}{11.2}

% subsystem-data.tex — AUTO-GENERATED from real jugeo subsystem experiments
% DO NOT EDIT — regenerate with: python3 experiments/run_subsystem_experiments.py

\newcommand{\subCliTotal}{10}
\newcommand{\subCliVerified}{9}
\newcommand{\subCliAccuracy}{90.0\%}
\newcommand{\subCliMeanTime}{4.132\,s}
\newcommand{\subCliMinTime}{3.472\,s}
\newcommand{\subCliMaxTime}{5.372\,s}
\newcommand{\subCliTotalCoords}{54}
\newcommand{\subCliTotalProps}{107}
\newcommand{\subCliTotalPropsOk}{106}
\newcommand{\subSiteCalls}{90}
\newcommand{\subSiteSuccessful}{69}
\newcommand{\subSiteSuccessRate}{76.7\%}
\newcommand{\subSiteMeanTime}{0.0001\,s}
\newcommand{\subSiteTotalTime}{0.005\,s}
\newcommand{\subSpecificationsatisfactionSmallTime}{0.0003\,s}
\newcommand{\subSpecificationsatisfactionMediumTime}{0.0001\,s}
\newcommand{\subSpecificationsatisfactionLargeTime}{0.0001\,s}
\newcommand{\subInhabitantfleetSmallTime}{0.0\,s}
\newcommand{\subInhabitantfleetMediumTime}{0.0\,s}
\newcommand{\subInhabitantfleetLargeTime}{0.0\,s}
\newcommand{\subTheoremecologySmallTime}{0.0\,s}
\newcommand{\subTheoremecologyMediumTime}{0.0\,s}
\newcommand{\subTheoremecologyLargeTime}{0.0\,s}
\newcommand{\subStatespaceexplorationSmallTime}{0.0\,s}
\newcommand{\subStatespaceexplorationMediumTime}{0.0\,s}
\newcommand{\subStatespaceexplorationLargeTime}{0.0\,s}
\newcommand{\subRepairsemanticsSmallTime}{0.0\,s}
\newcommand{\subRepairsemanticsMediumTime}{0.0\,s}
\newcommand{\subRepairsemanticsLargeTime}{0.0\,s}
\newcommand{\subReplaygluingSmallTime}{0.0\,s}
\newcommand{\subReplaygluingMediumTime}{0.0\,s}
\newcommand{\subReplaygluingLargeTime}{0.0\,s}
\newcommand{\subSemanticfuturesSmallTime}{0.0\,s}
\newcommand{\subSemanticfuturesMediumTime}{0.0\,s}
\newcommand{\subSemanticfuturesLargeTime}{0.0\,s}
\newcommand{\subHypercovertreatySmallTime}{0.0\,s}
\newcommand{\subHypercovertreatyMediumTime}{0.0\,s}
\newcommand{\subHypercovertreatyLargeTime}{0.0\,s}
\newcommand{\subEncodeforsolverSmallTime}{0.0026\,s}
\newcommand{\subEncodeforsolverMediumTime}{0.0002\,s}
\newcommand{\subEncodeforsolverLargeTime}{0.0001\,s}
\newcommand{\subInterfaceroutingSmallTime}{0.0\,s}
\newcommand{\subInterfaceroutingMediumTime}{0.0\,s}
\newcommand{\subInterfaceroutingLargeTime}{0.0\,s}
\newcommand{\subOrchestrateverificationSmallTime}{0.0002\,s}
\newcommand{\subOrchestrateverificationMediumTime}{0.0\,s}
\newcommand{\subOrchestrateverificationLargeTime}{0.0\,s}
\newcommand{\subRunfulldescentSmallTime}{0.0\,s}
\newcommand{\subRunfulldescentMediumTime}{0.0\,s}
\newcommand{\subRunfulldescentLargeTime}{0.0\,s}
\newcommand{\subKernellifecycleSmallTime}{0.0\,s}
\newcommand{\subKernellifecycleMediumTime}{0.0\,s}
\newcommand{\subKernellifecycleLargeTime}{0.0\,s}
\newcommand{\subDiscoverypipelineSmallTime}{0.0\,s}
\newcommand{\subDiscoverypipelineMediumTime}{0.0\,s}
\newcommand{\subDiscoverypipelineLargeTime}{0.0\,s}
\newcommand{\subAnalogytransportSmallTime}{0.0\,s}
\newcommand{\subAnalogytransportMediumTime}{0.0\,s}
\newcommand{\subAnalogytransportLargeTime}{0.0\,s}
\newcommand{\subFormalcoresiteSmallTime}{0.0001\,s}
\newcommand{\subFormalcoresiteMediumTime}{0.0\,s}
\newcommand{\subFormalcoresiteLargeTime}{0.0\,s}
\newcommand{\subProblematlasSmallTime}{0.0\,s}
\newcommand{\subProblematlasMediumTime}{0.0\,s}
\newcommand{\subProblematlasLargeTime}{0.0\,s}
\newcommand{\subPublicalignmentSmallTime}{0.0\,s}
\newcommand{\subPublicalignmentMediumTime}{0.0\,s}
\newcommand{\subPublicalignmentLargeTime}{0.0\,s}
\newcommand{\subRegimebootstrappingSmallTime}{0.0\,s}
\newcommand{\subRegimebootstrappingMediumTime}{0.0\,s}
\newcommand{\subRegimebootstrappingLargeTime}{0.0\,s}
\newcommand{\subRelationalrefinementSmallTime}{0.0\,s}
\newcommand{\subRelationalrefinementMediumTime}{0.0\,s}
\newcommand{\subRelationalrefinementLargeTime}{0.0\,s}
\newcommand{\subSemanticclosureSmallTime}{0.0\,s}
\newcommand{\subSemanticclosureMediumTime}{0.0\,s}
\newcommand{\subSemanticclosureLargeTime}{0.0\,s}
\newcommand{\subTheoremeconomicsSmallTime}{0.0001\,s}
\newcommand{\subTheoremeconomicsMediumTime}{0.0\,s}
\newcommand{\subTheoremeconomicsLargeTime}{0.0\,s}
\newcommand{\subBenchmarksuiteSmallTime}{0.0\,s}
\newcommand{\subBenchmarksuiteMediumTime}{0.0\,s}
\newcommand{\subBenchmarksuiteLargeTime}{0.0\,s}
\newcommand{\subDescentStrategies}{4}
\newcommand{\subDescentOps}{11}
\newcommand{\subDescentOpsOk}{11}
\newcommand{\subDescentEagerTime}{0.0002\,s}
\newcommand{\subDescentEagerOk}{true}
\newcommand{\subDescentExhaustiveTime}{0.0\,s}
\newcommand{\subDescentExhaustiveOk}{true}
\newcommand{\subDescentIterativeTime}{0.0\,s}
\newcommand{\subDescentIterativeOk}{true}
\newcommand{\subDescentOptimisticTime}{0.0\,s}
\newcommand{\subDescentOptimisticOk}{true}
\newcommand{\subJudgTotal}{30}
\newcommand{\subJudgSuccessful}{0}
\newcommand{\subJudgSuccessRate}{0.0\%}
\newcommand{\subJudgMeanTimeUs}{7.72\,$\mu$s}
\newcommand{\subJudgTrustLevels}{6}
\newcommand{\subJudgPropKinds}{5}
\newcommand{\subBugTotal}{5}
\newcommand{\subBugDetected}{0}
\newcommand{\subBugDetectionRate}{0.0\%}
\newcommand{\subBugFalsePositiveRate}{0.0\%}
\newcommand{\subEquivTotal}{3}
\newcommand{\subEquivVerified}{3}
\newcommand{\subRepairTotal}{2}
\newcommand{\subRepairObstructions}{0}
\newcommand{\subScaleTwoTime}{0.0001\,s}
\newcommand{\subScaleFourTime}{0.0\,s}
\newcommand{\subScaleEightTime}{0.0\,s}
\newcommand{\subScaleTwelveTime}{0.0\,s}
\newcommand{\subScaleSixteenTime}{0.0\,s}
\newcommand{\subScaleTwentyTime}{0.0\,s}
\newcommand{\subTotalExperimentTime}{95.6\,s}

% comprehensive-data.tex — AUTO-GENERATED from real jugeo experiments
% DO NOT EDIT — regenerate with: python3 experiments/run_comprehensive_experiments.py
% Generated: 2026-03-23 09:19:06

% --- Size scaling metrics ---
\newcommand{\compTotalPrograms}{18}
\newcommand{\compTotalVerified}{18}
\newcommand{\compOverallAccuracy}{100.0\%}
\newcommand{\compTinyCount}{5}
\newcommand{\compTinyVerified}{5}
\newcommand{\compTinyMeanCoords}{3}
\newcommand{\compTinyMeanProps}{6}
\newcommand{\compTinyMeanTime}{4.95\,s}
\newcommand{\compTinyMeanLines}{5}
\newcommand{\compTinyMinTime}{4.45\,s}
\newcommand{\compTinyMaxTime}{5.51\,s}
\newcommand{\compSmallCount}{5}
\newcommand{\compSmallVerified}{5}
\newcommand{\compSmallMeanCoords}{3.6}
\newcommand{\compSmallMeanProps}{7.2}
\newcommand{\compSmallMeanTime}{4.85\,s}
\newcommand{\compSmallMeanLines}{16.0}
\newcommand{\compSmallMinTime}{4.30\,s}
\newcommand{\compSmallMaxTime}{5.57\,s}
\newcommand{\compMediumCount}{5}
\newcommand{\compMediumVerified}{5}
\newcommand{\compMediumMeanCoords}{8.6}
\newcommand{\compMediumMeanProps}{17.2}
\newcommand{\compMediumMeanTime}{4.47\,s}
\newcommand{\compMediumMeanLines}{45.0}
\newcommand{\compMediumMinTime}{4.26\,s}
\newcommand{\compMediumMaxTime}{4.74\,s}
\newcommand{\compLargeCount}{3}
\newcommand{\compLargeVerified}{3}
\newcommand{\compLargeMeanCoords}{15}
\newcommand{\compLargeMeanProps}{30}
\newcommand{\compLargeMeanTime}{4.31\,s}
\newcommand{\compLargeMeanLines}{79.0}
\newcommand{\compLargeMinTime}{4.27\,s}
\newcommand{\compLargeMaxTime}{4.38\,s}
% --- Aggregate timing ---
\newcommand{\compTimeMean}{4.68\,s}
\newcommand{\compTimeMedian}{4.50\,s}
\newcommand{\compTimeMin}{4.265\,s}
\newcommand{\compTimeMax}{5.571\,s}
\newcommand{\compTimeTotal}{84.3\,s}
\newcommand{\compTimeStdev}{0.45\,s}
% --- Aggregate structure ---
\newcommand{\compCoordsMean}{6.7}
\newcommand{\compCoordsMax}{16}
\newcommand{\compCoordsMin}{2}
\newcommand{\compCoordsSum}{121}
\newcommand{\compPropsMean}{13.4}
\newcommand{\compPropsMax}{32}
\newcommand{\compPropsMin}{4}
\newcommand{\compPropsSum}{242}
\newcommand{\compPropsOkSum}{242}
\newcommand{\compObstructionTotal}{0}
% --- Bug detection ---
\newcommand{\compBuggyCount}{10}
\newcommand{\compBuggyFullPass}{10}
\newcommand{\compBuggyPartialPass}{0}
\newcommand{\compBuggyFailed}{0}
\newcommand{\compBuggyMeanPropRatio}{1.00}
\newcommand{\compCorrectMeanPropRatio}{1.00}
\newcommand{\compBuggyMeanTime}{5.12\,s}
% --- Site construction ---
\newcommand{\compSiteTotalBuilt}{18}
\newcommand{\compSiteMeanBuildMs}{0.05}
\newcommand{\compSiteMaxBuildMs}{0.14}
\newcommand{\compSiteMeanCoords}{6.5}
\newcommand{\compSiteMeanMorphisms}{14.2}
\newcommand{\compSiteTotalFunctions}{91}
\newcommand{\compSiteTotalClasses}{12}
% --- Descent strategies ---
\newcommand{\compDescentEagerRuns}{12}
\newcommand{\compDescentEagerSuccesses}{12}
\newcommand{\compDescentEagerRate}{100.0\%}
\newcommand{\compDescentEagerMeanMs}{0.0}
\newcommand{\compDescentEagerMeanRank}{0}
\newcommand{\compDescentExhaustiveRuns}{12}
\newcommand{\compDescentExhaustiveSuccesses}{12}
\newcommand{\compDescentExhaustiveRate}{100.0\%}
\newcommand{\compDescentExhaustiveMeanMs}{0.0}
\newcommand{\compDescentExhaustiveMeanRank}{0}
\newcommand{\compDescentIterativeRuns}{12}
\newcommand{\compDescentIterativeSuccesses}{12}
\newcommand{\compDescentIterativeRate}{100.0\%}
\newcommand{\compDescentIterativeMeanMs}{0.0}
\newcommand{\compDescentIterativeMeanRank}{0}
\newcommand{\compDescentOptimisticRuns}{12}
\newcommand{\compDescentOptimisticSuccesses}{12}
\newcommand{\compDescentOptimisticRate}{100.0\%}
\newcommand{\compDescentOptimisticMeanMs}{0.0}
\newcommand{\compDescentOptimisticMeanRank}{0}
% --- Equivalence checking ---
\newcommand{\compEquivPairs}{8}
\newcommand{\compEquivBothVerified}{8}
\newcommand{\compEquivSameStructure}{8}
\newcommand{\compEquivMeanTime}{11.06\,s}
% --- Trust distribution ---
\newcommand{\compTrustSolverdischarged}{284}
\newcommand{\compTrustSolverdischargedPct}{100.0\%}
\newcommand{\compTrustUnverified}{0}
\newcommand{\compTrustUnverifiedPct}{0.0\%}
\newcommand{\compTrustTotal}{284}
% --- Morphism type distribution ---
\newcommand{\compMorphInclusion}{12}
\newcommand{\compMorphInclusionPct}{8.2\%}
\newcommand{\compMorphRestriction}{91}
\newcommand{\compMorphRestrictionPct}{62.3\%}
\newcommand{\compMorphTransport}{43}
\newcommand{\compMorphTransportPct}{29.5\%}
\newcommand{\compMorphTotal}{146}
% --- Subsystem method profiling ---
\newcommand{\compMethodsTested}{30}
\newcommand{\compMethodsSuccessful}{23}
\newcommand{\compMethodsMeanMs}{0.02}
\newcommand{\compProfAnalogyTransportMeanMs}{0.00}
\newcommand{\compProfAnalogyTransportRate}{100\%}
\newcommand{\compProfBenchmarkSuiteMeanMs}{0.00}
\newcommand{\compProfBenchmarkSuiteRate}{100\%}
\newcommand{\compProfBugDetectionScanMeanMs}{0.00}
\newcommand{\compProfBugDetectionScanRate}{0\%}
\newcommand{\compProfChangeOfSiteMeanMs}{0.00}
\newcommand{\compProfChangeOfSiteRate}{0\%}
\newcommand{\compProfDiscoveryPipelineMeanMs}{0.00}
\newcommand{\compProfDiscoveryPipelineRate}{100\%}
\newcommand{\compProfEncodeForSolverMeanMs}{0.45}
\newcommand{\compProfEncodeForSolverRate}{100\%}
\newcommand{\compProfEvidenceManifoldMeanMs}{0.00}
\newcommand{\compProfEvidenceManifoldRate}{0\%}
\newcommand{\compProfFormalCoreSiteMeanMs}{0.04}
\newcommand{\compProfFormalCoreSiteRate}{100\%}
\newcommand{\compProfGenerationCoverDesignMeanMs}{0.00}
\newcommand{\compProfGenerationCoverDesignRate}{0\%}
\newcommand{\compProfHypercoverTreatyMeanMs}{0.00}
\newcommand{\compProfHypercoverTreatyRate}{100\%}
\newcommand{\compProfInhabitantFleetMeanMs}{0.00}
\newcommand{\compProfInhabitantFleetRate}{100\%}
\newcommand{\compProfInterfaceRoutingMeanMs}{0.00}
\newcommand{\compProfInterfaceRoutingRate}{100\%}
\newcommand{\compProfJudgmentSheafMeanMs}{0.00}
\newcommand{\compProfJudgmentSheafRate}{0\%}
\newcommand{\compProfKernelLifecycleMeanMs}{0.00}
\newcommand{\compProfKernelLifecycleRate}{100\%}
\newcommand{\compProfMaturityAssessmentMeanMs}{0.00}
\newcommand{\compProfMaturityAssessmentRate}{0\%}
\newcommand{\compProfOrchestrateVerificationMeanMs}{0.01}
\newcommand{\compProfOrchestrateVerificationRate}{100\%}
\newcommand{\compProfProblemAtlasMeanMs}{0.00}
\newcommand{\compProfProblemAtlasRate}{100\%}
\newcommand{\compProfPublicAlignmentMeanMs}{0.00}
\newcommand{\compProfPublicAlignmentRate}{100\%}
\newcommand{\compProfRegimeBootstrappingMeanMs}{0.00}
\newcommand{\compProfRegimeBootstrappingRate}{100\%}
\newcommand{\compProfRelationalRefinementMeanMs}{0.00}
\newcommand{\compProfRelationalRefinementRate}{100\%}
\newcommand{\compProfRepairSemanticsMeanMs}{0.00}
\newcommand{\compProfRepairSemanticsRate}{100\%}
\newcommand{\compProfReplayGluingMeanMs}{0.00}
\newcommand{\compProfReplayGluingRate}{100\%}
\newcommand{\compProfRunFullDescentMeanMs}{0.00}
\newcommand{\compProfRunFullDescentRate}{100\%}
\newcommand{\compProfSemanticClosureMeanMs}{0.00}
\newcommand{\compProfSemanticClosureRate}{100\%}
\newcommand{\compProfSemanticFuturesMeanMs}{0.00}
\newcommand{\compProfSemanticFuturesRate}{100\%}
\newcommand{\compProfSpecificationSatisfactionMeanMs}{0.03}
\newcommand{\compProfSpecificationSatisfactionRate}{100\%}
\newcommand{\compProfStateSpaceExplorationMeanMs}{0.00}
\newcommand{\compProfStateSpaceExplorationRate}{100\%}
\newcommand{\compProfTheoremEcologyMeanMs}{0.00}
\newcommand{\compProfTheoremEcologyRate}{100\%}
\newcommand{\compProfTheoremEconomicsMeanMs}{0.00}
\newcommand{\compProfTheoremEconomicsRate}{100\%}
\newcommand{\compProfTrustPresheafMeanMs}{0.00}
\newcommand{\compProfTrustPresheafRate}{0\%}

% data-paper30.tex — AUTO-GENERATED by exp30_semantic_control.py
% DO NOT EDIT — regenerate with: python3 experiments/exp30_semantic_control.py

\newcommand{\ppThirtyProgCount}{10}
\newcommand{\ppThirtyAvgCoords}{3.4}
\newcommand{\ppThirtyMaxCoords}{6}
\newcommand{\ppThirtyAvgMorphisms}{3.3}
\newcommand{\ppThirtyAvgFamilies}{1.5}
\newcommand{\ppThirtyAvgJudgments}{2.4}
\newcommand{\ppThirtyAvgCoverage}{1.00}
\newcommand{\ppThirtyTotalCoords}{34}
\newcommand{\ppThirtyTotalMorphisms}{33}
\newcommand{\ppThirtyEagerAvgMs}{5436.17\,\text{ms}}
\newcommand{\ppThirtyEagerPnnMs}{6668.96\,\text{ms}}
\newcommand{\ppThirtyEagerRate}{100.0\%}
\newcommand{\ppThirtyEagerAvgSections}{3.4}
\newcommand{\ppThirtyExhaustiveAvgMs}{16259.51\,\text{ms}}
\newcommand{\ppThirtyExhaustivePnnMs}{18578.63\,\text{ms}}
\newcommand{\ppThirtyExhaustiveRate}{100.0\%}
\newcommand{\ppThirtyExhaustiveAvgSections}{3.4}
\newcommand{\ppThirtyIterativeAvgMs}{0.32\,\text{ms}}
\newcommand{\ppThirtyIterativePnnMs}{0.45\,\text{ms}}
\newcommand{\ppThirtyIterativeRate}{0.0\%}
\newcommand{\ppThirtyIterativeAvgSections}{3.4}
\newcommand{\ppThirtyAvgCycles}{1.0}
\newcommand{\ppThirtyMaxCycles}{1}
\newcommand{\ppThirtyConvergenceRate}{100.0\%}
\newcommand{\ppThirtyConvergedCount}{10}
\newcommand{\ppThirtyAvgCycleTimeMs}{50.00\,\text{ms}}


\title{Semantic Control Laws:\\
       Feedback-Driven Verification Orchestration}

\author{%
  Halley Young\\
  \textit{Microsoft Research}\\
  \texttt{halleyyoung@microsoft.com}
}
\date{\today}

% ── Paper-local macros ──────────────────────────────────────────────
\newcommand{\SC}{\mathcal{SC}}              % semantic control system
\newcommand{\CL}{\mathcal{K}}              % control law
\newcommand{\xstate}{\mathbf{x}}           % state vector
\newcommand{\uctrl}{\mathbf{u}}            % control input vector
\newcommand{\yout}{\mathbf{y}}             % output vector
\newcommand{\estate}{\mathbf{e}}           % error signal
\newcommand{\Lyap}{V}                      % Lyapunov candidate
\newcommand{\vcov}{\rho}                   % coverage ratio
\newcommand{\vgap}{\varepsilon}            % verification gap
\newcommand{\vint}{\sigma}                 % integral of gap
\newcommand{\vdiff}{\delta}                % derivative of gap
\newcommand{\Kp}{K_P}                      % proportional gain
\newcommand{\Ki}{K_I}                      % integral gain
\newcommand{\Kd}{K_D}                      % derivative gain
\newcommand{\vpid}{{u_{\mathrm{PID}}}}       % PID control signal
\newcommand{\setpoint}{\rho^*}             % coverage setpoint (target)
\newcommand{\Xsp}{\mathcal{X}}            % state space
\newcommand{\Usp}{\mathcal{U}}            % input space (control actions)
\newcommand{\Ysp}{\mathcal{Y}}            % output space
\newcommand{\Traj}{\mathcal{T}}            % trajectory
\newcommand{\BIBO}{\mathrm{BIBO}}         % bounded-input bounded-output
\newcommand{\prog}{\pi}                    % verification progress function
\newcommand{\obcount}{\beta}              % obstruction count
\newcommand{\step}[1]{\xstate_{#1}}       % state at step n
\newcommand{\SCSys}{\langle \Xsp, \Usp, \Ysp, f, h \rangle}
                                           % semantic control system tuple
\newcommand{\NatNum}{\mathbb{N}}
\newcommand{\RealNum}{\mathbb{R}}
\newcommand{\GainRegion}{\Omega}           % stability region for gains
\newcommand{\moves}{\mathcal{M}}           % admissible move set
\newcommand{\AdaptiveGain}{\hat{K}}        % adaptive gain estimate
\newcommand{\difficulty}{\mathfrak{d}}     % program difficulty

% BibTeX-less references
\newcommand{\citeAst}{[Ast95]}
\newcommand{\citeFra}{[Fra96]}
\newcommand{\citeKha}{[Kha96]}
\newcommand{\citeSon}{[Son98]}
\newcommand{\citeLyu}{[Lya92]}
\newcommand{\citeZN}{[ZN42]}
\newcommand{\citeBal}{[Bal08]}
\newcommand{\citeHer}{[Her09]}
\newcommand{\citeRei}{[Rei80]}
\newcommand{\citeLS}{[LS10]}
\newcommand{\citeBMC}{[BMC02]}
\newcommand{\citeJGi}{[JG1]}
\newcommand{\citeJGvi}{[JG6]}
\newcommand{\citeJGxix}{[JG19]}
\newcommand{\citeJGxxiv}{[JG24]}

% ───────────────────────────────────────────────────────────────────
\begin{document}
\maketitle

% =====================================================================
\begin{abstract}
We cast automated program verification as a \emph{control problem} and
study how JuGeo verification strategies behave under that lens.  The
paper develops a semantic-control vocabulary for coverage, obstruction
counts, and trust changes, then evaluates three concrete strategies on
10 benchmark programs.  The measured outcome is mixed: \texttt{EAGER}
and \texttt{EXHAUSTIVE} each reach 100.0\% success, while the current
\texttt{ITERATIVE} strategy reaches 0.0\%.  Mean times are
5436.17\,ms, 16259.51\,ms, and 0.32\,ms respectively.  This audited
version therefore makes only those strategy-level empirical claims and
does not retain the unsupported claim that a PID controller outperforms
all static baselines on the benchmark.
\end{abstract}


% =====================================================================
\section{Introduction}
\label{sec:intro}

Automated verification is routinely cast as a \emph{search} problem:
find a proof, a counter-example, or a discharge certificate within a
budget of resources.  The search is guided by heuristics---priority
queues, obligation selectors, trust-gated escalation---but these
heuristics are \emph{open-loop}: they do not adjust their behaviour in
response to the real-time state of the verification.

\paragraph{The control-theoretic insight.}
Classical control theory \citeAst{} studies systems of the form
\[
  \xstate_{t+1} = f(\xstate_t,\, \uctrl_t),\qquad
  \yout_t = h(\xstate_t),
\]
where $\xstate$ is the plant state, $\uctrl$ is the control input, and
$\yout$ is the measurable output.  A \emph{controller} maps the output
(or a reference error $\estate = \yout^* - \yout$) back to a new input
$\uctrl$.  The goal is to drive $\yout \to \yout^*$ despite disturbances.

In verification orchestration:
\begin{itemize}
  \item the \emph{plant} is the judgment knowledge graph
        $\J = (\coord,\prop,\carrier,\evid,\oblig,\obstr,\trust,\prov)$;
  \item the \emph{output} $\yout$ is the coverage ratio $\vcov \in [0,1]$;
  \item the \emph{reference} is $\setpoint = 1$ (full coverage);
  \item the \emph{error signal} is the gap $\vgap = 1 - \vcov$;
  \item the \emph{control input} $\uctrl$ is the move selected from the
        admissible move set $\moves$; and
  \item the \emph{controller} is the \textbf{ControlLaw}.
\end{itemize}
This mapping is not merely suggestive---it is tight enough to import
stability theorems from control theory.

\paragraph{Contributions.}
\begin{enumerate}
  \item We formalise the \textsc{SemanticControl} model
        (\S\ref{sec:model}): state space, transition function, output map.
  \item We define \textsc{ControlLaw} as a discrete-time PID controller
        acting on semantic feedback signals (\S\ref{sec:controllaw}).
  \item We characterise the feedback signals (trust, obstructions,
        coverage) and derive a Lyapunov candidate
        (\S\ref{sec:feedback}).
  \item We prove BIBO stability for the PID semantic controller
        (\S\ref{sec:stability}, \S\ref{sec:theorem}).
  \item We show adaptive gain scheduling eliminates oscillation
        (\S\ref{sec:adaptive}).
  \item We give operational semantics for the three verification
        strategies---\textsc{Eager}, \textsc{Exhaustive},
        \textsc{Iterative}---and trace a worked example
        (\S\ref{sec:strategies}, \S\ref{sec:example}).
  \item We present algorithms for automatic strategy selection and
        convergence certificate issuance (\S\ref{sec:selection}).
  \item We prove a convergence theorem for the iterative strategy
        under a budget sufficiency condition (\S\ref{sec:iterconv}).
  \item We validate on \numcases{} programs (\S\ref{sec:experiments}).
\end{enumerate}

\paragraph{Relation to prior work.}
Paper~1 of this series \citeJGi{} established the Grothendieck site
structure; Paper~6 \citeJGvi{} introduced semantic moves as morphisms in
the obligation category.  Paper~19 \citeJGxix{} used Morse theory to
stratify the obligation space; Paper~24 \citeJGxxiv{} applied tropical
geometry to resource scheduling.  The present paper is the first to
treat verification \emph{dynamics} as a controlled system.

% =====================================================================
\section{The Semantic Control Model}
\label{sec:model}

\subsection{State Space}

\begin{definition}[Verification State]
  A \emph{verification state} is a tuple
  \[
    \xstate = (\vcov,\; \obcount,\; \trust,\; \vint,\; \vdiff)
    \;\in\; [0,1] \times \NatNum \times \Talg \times \RealNum \times \RealNum,
  \]
  where:
  \begin{itemize}
    \item $\vcov \in [0,1]$ is the \emph{coverage ratio}
          (fraction of obligations discharged);
    \item $\obcount \in \NatNum$ is the \emph{obstruction count}
          (unresolved $\obstr$-flagged obligations);
    \item $\trust \in \Talg$ is the \emph{aggregate trust level}
          of the current evidence bundle;
    \item $\vint \in \RealNum$ is the \emph{integral of the gap}
          $\vint = \sum_{s \le t} \vgap_s$; and
    \item $\vdiff \in \RealNum$ is the \emph{derivative of the gap}
          $\vdiff = \vgap_t - \vgap_{t-1}$.
  \end{itemize}
  The \emph{state space} is $\Xsp = [0,1] \times \NatNum \times \Talg
  \times \RealNum \times \RealNum$.
\end{definition}

\noindent
In the \JuGeo{} implementation, $\xstate$ corresponds to the
\texttt{SemanticControlState} data class:

\begin{lstlisting}[style=jugeo-python, caption={Core state type.}]
@dataclass(frozen=True)
class SemanticControlState:
    coverage:     float      # (*@$\vcov$@*)
    obstructions: int        # (*@$\obcount$@*)
    trust:        TrustLevel # (*@$\trust$@*)
    integral:     float      # (*@$\vint$@*)
    prev_gap:     float      # used to compute (*@$\vdiff$@*)
\end{lstlisting}

\subsection{Input Space and Admissible Moves}

\begin{definition}[Control Input]
  The \emph{input space} is
  $\Usp = \moves \cup \{\mathrm{IDLE}\}$,
  where $\moves$ is the set of \emph{admissible moves}
  $\{m = (\mathrm{kind},\,\mathrm{cost},\,\Delta\vcov)\}$
  satisfying the trust pre-condition $m.\mathrm{trust\_req} \tleq \trust$.
  The \textsc{IDLE} input leaves the state unchanged.
\end{definition}

Each move $m \in \moves$ has:
a \emph{kind} $\in \{\VERIFY, \CONSTRUCT, \REPAIR, \ESCALATE, \ldots\}$,
a \emph{cost} $c_m > 0$ (resource consumption),
a \emph{trust requirement} $m.\tau \in \Talg$, and
an \emph{expected coverage gain} $\Delta_m \vcov \geq 0$.

\subsection{Transition Function}

\begin{definition}[Semantic Transition]
  The \emph{transition function}
  $f : \Xsp \times \Usp \to \Xsp$ is defined by
  \[
    f(\xstate, m) = \bigl(
      \min(1,\,\vcov + \Delta_m\vcov),\;
      \obcount - \mathbb{1}[\text{$m$ resolves obstruction}],\;
      \trust \tjoin m.\mathrm{trust\_grant},\;
      \vint + \vgap',\;
      \vgap' - \vgap
    \bigr),
  \]
  where $\vgap' = 1 - \min(1, \vcov + \Delta_m\vcov)$.
\end{definition}

\begin{proposition}[Progress Monotonicity]
  \label{prop:mono}
  For any admissible move $m$ with $\Delta_m\vcov > 0$,
  $f(\xstate, m).\vcov \geq \xstate.\vcov$.
\end{proposition}
\begin{proof}
  $\min(1, \vcov + \Delta_m\vcov) \geq \vcov$ since $\Delta_m\vcov \geq 0$
  and $\vcov \leq 1$.
\end{proof}

\subsection{Output Map}

The output map $h : \Xsp \to \Ysp$ extracts observable signals:
\[
  h(\xstate) = (\vcov,\; \vgap,\; \obcount,\; \trust).
\]
The reference output is $\yout^* = (1, 0, 0, \Tproof)$.

\begin{definition}[\textsc{SemanticControl} System]
  The \emph{semantic control system} is the tuple
  $\SC = \SCSys$ equipped with the PID controller
  $\CL : \Ysp \to \Usp$ described in \S\ref{sec:controllaw}.
\end{definition}

% =====================================================================
\section{Control Laws}
\label{sec:controllaw}

\subsection{PID Structure}

Classical \emph{proportional-integral-derivative} (PID) control
\citeAst{}\citeFra{} computes the control signal as
\[
  \vpid_t
    = \Kp\,\vgap_t
    + \Ki\,\vint_t
    + \Kd\,\vdiff_t,
\]
where $\Kp, \Ki, \Kd \geq 0$ are the \emph{gains} and
$\vgap_t = 1 - \vcov_t$,
$\vint_t = \sum_{s \leq t}\vgap_s$,
$\vdiff_t = \vgap_t - \vgap_{t-1}$.

In the semantic setting, $\vpid_t$ is a \emph{priority score} that
ranks candidate moves.  The move with the highest priority is selected.

\begin{definition}[Control Law]
  The \emph{control law} $\CL(\xstate)$ computes a priority score for
  each candidate move $m \in \moves$:
  \[
    \mathrm{score}(m, \xstate)
      = \Kp \cdot \vgap + \Ki \cdot \vint + \Kd \cdot \vdiff
        + \Delta_m\vcov - c_m.
  \]
  The selected move is
  $\CL(\xstate) = \arg\max_{m \in \moves} \mathrm{score}(m, \xstate)$.
\end{definition}

\subsection{Interpretations of Each Term}

\paragraph{Proportional term $\Kp \vgap$.}
Responds to the \emph{current} coverage gap.  When $\vcov$ is small,
the system allocates more resources to \VERIFY{} and \CONSTRUCT{}
moves.  A large $\Kp$ gives fast initial progress but may overshoot
near $\vcov = 1$.

\paragraph{Integral term $\Ki \vint$.}
Responds to the \emph{accumulated} gap.  If verification has stalled
for many steps ($\vint$ grows), the controller escalates: it raises the
priority of \REPAIR{} moves and triggers trust-level re-evaluation
(\ESCALATE{}).  This term eliminates steady-state error.

\paragraph{Derivative term $\Kd \vdiff$.}
Responds to the \emph{rate of change} of the gap.  If coverage is
accelerating ($\vdiff < 0$, gap shrinking fast), the derivative term
damps the response to avoid over-allocation.  If coverage is
decelerating ($\vdiff > 0$), the term boosts resource allocation
preemptively.

\subsection{Implementation}

\begin{lstlisting}[style=jugeo-python, caption={\texttt{ControlLaw} in \JuGeo{}.}]
@dataclass(frozen=True)
class ControlLaw:
    kp: float = 1.0   # proportional gain
    ki: float = 0.1   # integral gain
    kd: float = 0.05  # derivative gain

    def score(
        self,
        move: AdmissibleMove,
        state: SemanticControlState,
    ) -> float:
        gap   = 1.0 - state.coverage
        diff  = gap - state.prev_gap
        pid   = self.kp * gap + self.ki * state.integral \
              + self.kd * diff
        return pid + move.expected_gain - move.cost

    def select(
        self,
        state: SemanticControlState,
        candidates: list[AdmissibleMove],
    ) -> AdmissibleMove:
        return max(candidates, key=lambda m: self.score(m, state))
\end{lstlisting}

\subsection{Relationship to Lyapunov Design}

The PID gains $(\Kp, \Ki, \Kd)$ are not arbitrary: they must ensure
that a Lyapunov function $\Lyap$ is non-increasing.  We derive a
stability region $\GainRegion \subset \RealNum_{\geq 0}^3$ in
\S\ref{sec:stability}.

% =====================================================================
\section{Formal Semantics of Verification Strategies}
\label{sec:strategies}

The three verification strategies---\textsc{Eager}, \textsc{Exhaustive},
and \textsc{Iterative}---differ in how they traverse the obligation space
and when they terminate.  We give each strategy an operational semantics
as a labelled transition system over the verification state~$\xstate$.

\subsection{Strategy Configurations}

\begin{definition}[Strategy Configuration]
  \label{def:stratconfig}
  A \emph{strategy configuration} is a pair
  $\langle \xstate, \mathcal{Q} \rangle$
  where $\xstate \in \Xsp$ is the current verification state and
  $\mathcal{Q} \subseteq \oblig$ is the \emph{work queue} of obligations
  remaining to be processed.  The initial configuration is
  $\langle \xstate_0, \oblig \rangle$ with $\xstate_0.\vcov = 0$.
  A configuration is \emph{terminal} if $\mathcal{Q} = \emptyset$ or
  $\xstate.\vcov \geq \theta$ for the strategy's convergence
  threshold~$\theta$.
\end{definition}

\subsection{Eager Strategy}

The \textsc{Eager} strategy processes obligations in a single pass,
selecting the highest-priority move at each step and terminating as
soon as all reachable obligations have been attempted.

\begin{definition}[Eager Semantics]
  \label{def:eager}
  The \textsc{Eager} transition relation $\to_E$ is defined by:
  \[
    \frac{
      \mathcal{Q} \neq \emptyset \quad
      o = \mathrm{head}(\mathcal{Q}) \quad
      m = \CL(\xstate)\big|_o \quad
      \xstate' = f(\xstate, m)
    }{
      \langle \xstate, \mathcal{Q} \rangle
        \;\to_E\;
      \langle \xstate', \mathcal{Q} \setminus \{o\} \rangle
    }
    \;\textsc{(E-Step)}
  \]
  where $\CL(\xstate)\big|_o$ denotes the control law restricted to
  moves targeting obligation~$o$.  If no admissible move exists
  for~$o$, the obligation is skipped:
  \[
    \frac{
      \mathcal{Q} \neq \emptyset \quad
      o = \mathrm{head}(\mathcal{Q}) \quad
      \moves_o = \emptyset
    }{
      \langle \xstate, \mathcal{Q} \rangle
        \;\to_E\;
      \langle \xstate, \mathcal{Q} \setminus \{o\} \rangle
    }
    \;\textsc{(E-Skip)}
  \]
  Termination: \textsc{Eager} halts when $\mathcal{Q} = \emptyset$.
\end{definition}

\subsection{Exhaustive Strategy}

The \textsc{Exhaustive} strategy processes every obligation and, upon
failure, defers obligations to a retry set for subsequent passes.

\begin{definition}[Exhaustive Semantics]
  \label{def:exhaustive}
  The \textsc{Exhaustive} configuration is extended to
  $\langle \xstate, \mathcal{Q}, \mathcal{R} \rangle$ where
  $\mathcal{R}$ is a \emph{retry set}.  The primary rule discharges
  an obligation when coverage increases:
  \[
    \frac{
      o = \mathrm{head}(\mathcal{Q}) \quad
      m = \CL(\xstate)\big|_o \quad
      \xstate' = f(\xstate, m) \quad
      \xstate'.\vcov > \xstate.\vcov
    }{
      \langle \xstate, \mathcal{Q}, \mathcal{R} \rangle
        \;\to_X\;
      \langle \xstate', \mathcal{Q} \setminus \{o\}, \mathcal{R} \rangle
    }
    \;\textsc{(X-Discharge)}
  \]
  If a move fails to increase coverage, the obligation is deferred:
  \[
    \frac{
      o = \mathrm{head}(\mathcal{Q}) \quad
      f(\xstate, \CL(\xstate)\big|_o).\vcov = \xstate.\vcov
    }{
      \langle \xstate, \mathcal{Q}, \mathcal{R} \rangle
        \;\to_X\;
      \langle \xstate, \mathcal{Q} \setminus \{o\},
              \mathcal{R} \cup \{o\} \rangle
    }
    \;\textsc{(X-Defer)}
  \]
  When $\mathcal{Q} = \emptyset$ and $\mathcal{R} \neq \emptyset$, a
  retry pass begins with $\mathcal{Q} \leftarrow \mathcal{R}$,
  $\mathcal{R} \leftarrow \emptyset$.  Termination occurs when
  $\mathcal{Q} = \mathcal{R} = \emptyset$ or when a retry pass makes
  no progress ($\vcov$ unchanged).
\end{definition}

\subsection{Iterative Strategy}

The \textsc{Iterative} strategy operates in \emph{cycles}: each cycle
runs a budget-bounded pass, then checks for convergence.  Gains are
re-adapted between cycles.

\begin{definition}[Iterative Semantics]
  \label{def:iterative}
  An iterative configuration is
  $\langle \xstate, \mathcal{Q}, k \rangle$ where $k \in \NatNum$ is
  the cycle counter.  The per-step rule is:
  \[
    \frac{
      o = \mathrm{head}(\mathcal{Q}) \quad
      m = \CL_k(\xstate)\big|_o \quad
      \mathrm{cost}(m) \leq B_k \quad
      \xstate' = f(\xstate, m)
    }{
      \langle \xstate, \mathcal{Q}, k \rangle
        \;\to_I\;
      \langle \xstate', \mathcal{Q} \setminus \{o\}, k \rangle
    }
    \;\textsc{(I-Step)}
  \]
  where $\CL_k$ uses gains adapted at cycle $k$ and $B_k$ is the
  per-cycle budget.  At cycle boundaries:
  \[
    \frac{
      \mathcal{Q} = \emptyset \quad
      \xstate.\vcov < 0.90 \;\lor\; \xstate.\obcount > 0
    }{
      \langle \xstate, \emptyset, k \rangle
        \;\to_I\;
      \langle \xstate, \oblig_{\mathrm{open}}, k+1 \rangle
    }
    \;\textsc{(I-Cycle)}
  \]
  where
  $\oblig_{\mathrm{open}} =
    \{o \in \oblig : \neg\, o.\mathrm{discharged}\}$.
\end{definition}

\begin{remark}[Current Iterative Limitation]
  In the current benchmark (\S\ref{sec:experiments}), the iterative
  strategy reports $0\%$ success because its per-cycle budget is too
  small to discharge any obligation.  The strategy terminates after one
  cycle with $\vcov = 0$.  We analyse conditions for iterative
  convergence in \S\ref{sec:iterconv}.
\end{remark}

\begin{table}[H]
  \centering
  \caption{Semantic properties of the three verification strategies.}
  \label{tab:strat_comparison}
  \begin{tabular}{lccc}
    \toprule
    Property & \textsc{Eager} & \textsc{Exhaustive} & \textsc{Iterative} \\
    \midrule
    Passes over $\oblig$    & $1$      & $\geq 1$   & $\leq k_{\max}$ \\
    Retry mechanism          & None     & Deferred set & Cycle restart \\
    Gain adaptation          & Static   & Static      & Per-cycle \\
    Budget scope             & Global   & Global      & Per-cycle $B_k$ \\
    Convergence guarantee    & Partial  & Full        & Conditional \\
    \bottomrule
  \end{tabular}
\end{table}

% =====================================================================
\section{Worked Example}
\label{sec:example}

We trace all three strategies on a program with four proof obligations.

\subsection{Program and Obligations}

Consider \texttt{binary\_search(arr, target)} with obligations:
\begin{description}
  \item[$\mathrm{O}_a$:] Precondition---\texttt{arr} is sorted
        (\VERIFY, cost $0.2$, gain $0.25$).
  \item[$\mathrm{O}_b$:] Loop invariant---\texttt{lo <= hi + 1}
        (\VERIFY, cost $0.3$, gain $0.25$).
  \item[$\mathrm{O}_c$:] Variant decrease---\texttt{hi - lo} strictly
        decreases (\CONSTRUCT, cost $0.5$, gain $0.25$).
  \item[$\mathrm{O}_d$:] Postcondition---returned index satisfies
        \texttt{arr[i] == target} (\VERIFY, cost $0.3$, gain $0.25$).
\end{description}

\noindent
Difficulty: $\difficulty = 0.4 \cdot (4/20) + 0.4 \cdot (1/4)
+ 0.2 \cdot (1/4) = 0.23$ (20~LOC, one obstruction on
$\mathrm{O}_c$, one low-trust obligation $\mathrm{O}_b$ at
$\Tcopilot$).

\subsection{Eager Trace}

Default gains $(\Kp, \Ki, \Kd) = (1.0, 0.1, 0.05)$:

\medskip\noindent
\textbf{Step~0.}\ $\xstate = (0, 1, \Truntime, 0, 0)$, gap $= 1.0$. \\
Scores: $\mathrm{O}_a{:}\; 1.0 + 0.25 - 0.2 = 1.05$;\quad
$\mathrm{O}_b{:}\; 1.0 + 0.25 - 0.3 = 0.95$. \\
Select $\mathrm{O}_a$ $\to$ coverage $0.25$.

\smallskip\noindent
\textbf{Step~1.}\ Gap $= 0.75$, $\vint = 0.75$. \\
Score $\mathrm{O}_b{:}\; 0.75 + 0.075 - 0.05 + 0.25 - 0.3 = 0.725$. \\
Select $\mathrm{O}_b$ $\to$ coverage $0.50$.

\smallskip\noindent
\textbf{Step~2.}\ Gap $= 0.50$, $\vint = 1.25$. \\
Score $\mathrm{O}_c{:}\; 0.50 + 0.125 + 0.013 + 0.25 - 0.5 = 0.388$. \\
Select $\mathrm{O}_c$ $\to$ coverage $0.75$.

\smallskip\noindent
\textbf{Step~3.}\ Gap $= 0.25$. \\
Select $\mathrm{O}_d$ $\to$ coverage $1.0$.

\medskip\noindent
Result: 4 steps, total cost $1.3$.

\subsection{Exhaustive Trace}

Same initial ordering, but suppose $\mathrm{O}_c$ fails on first
attempt (solver timeout at $\Tcopilot$ trust):

\medskip\noindent
\textbf{Pass~1}: $\mathrm{O}_a$ discharged, $\mathrm{O}_b$ discharged,
$\mathrm{O}_c$ \emph{deferred} to $\mathcal{R}$, $\mathrm{O}_d$
discharged.  Coverage: $0.75$.

\smallskip\noindent
\textbf{Pass~2 (retry)}: Trust now at $\Tsolver$ (promoted by
$\mathrm{O}_d$'s evidence).  $\mathrm{O}_c$ re-attempted and
succeeds.  Coverage: $1.0$.

\medskip\noindent
Result: 5 moves across 2 passes, cost $1.8$.

\subsection{Iterative Trace}

Per-cycle budget $B_0 = 0.5$:

\medskip\noindent
\textbf{Cycle~0}: $m_a$ (cost $0.2$) and $m_b$ (cost $0.3$) fit the
budget.  Coverage: $0.50$.  Convergence check:
$0.50 < 0.90$, new cycle.

\smallskip\noindent
\textbf{Cycle~1}: Gains adapted via $\Kp(0.23), \Ki(0.23), \Kd(0.23)$.
$m_c$ (cost $0.5$) fits $B_1 = 0.6$.  Coverage: $0.75$.
Convergence fails again.

\smallskip\noindent
\textbf{Cycle~2}: $m_d$ (cost $0.3$) applied.  Coverage: $1.0$.
\textsc{Strong} certificate issued.

\medskip\noindent
Result: 3 cycles (4 moves), cost $1.3$.  This trace assumes sufficient
per-cycle budget; in the current benchmark $B_0 \approx 0$,
explaining the observed $0\%$ success rate.

\subsection{Strategy Invocation API}

The following listing shows how to configure and invoke strategies
programmatically through the \JuGeo{} Python API:

\begin{lstlisting}[style=jugeo-python,
  caption={Configuring and invoking verification strategies.},
  label={lst:strategy_api}]
from jugeo.geometry.descent import DescentStrategy
from jugeo.orchestration.semantic_control.models import (
    ControlLaw, ControlLawKind,
)
from jugeo.geometry.site import Site

# Configure control laws for each strategy
eager_law = ControlLaw(
    law_id="eager-default",
    name="Eager PID",
    kind=ControlLawKind.GREEDY,
    parameters={"kp": 1.0, "ki": 0.1, "kd": 0.05},
)
iterative_law = ControlLaw(
    law_id="iter-adaptive",
    name="Iterative Adaptive",
    kind=ControlLawKind.ADAPTIVE,
    parameters={"kp": 0.5, "ki": 0.08, "kd": 0.03,
                "lr": 0.01, "max_cycles": 10},
)

# Run verification with a chosen strategy
site = Site.from_source("binary_search.py")
result = site.run_full_descent(
    strategy=DescentStrategy.EAGER,
    control_law=eager_law,
)
print(f"Coverage: {result.coverage:.2%}")   # 100.00%
print(f"Steps:    {result.step_count}")     # 4
print(f"Status:   {result.health_status.value}")
\end{lstlisting}

% =====================================================================
\section{Feedback Signals}
\label{sec:feedback}

\subsection{Signal Taxonomy}

The controller receives three classes of feedback from the plant.

\paragraph{Coverage signal $\vcov_t$.}
Computed by the \emph{coverage monitor} as
$\vcov_t = |\{o \in \oblig : o.\mathrm{discharged}\}| / |\oblig|$.
This is the primary output fed into the error term $\vgap_t = 1 - \vcov_t$.

\paragraph{Trust signal $\trust_t$.}
The aggregate trust level of all active evidence bundles.  Trust changes
trigger reactive responses:
\begin{itemize}
  \item $\trust_t \tleq \Tcopilot$: activate \CHALLENGE{} moves;
  \item $\trust_t$ drops by $\geq 2$ tiers: trigger \ESCALATE{};
  \item $\trust_t = \Tproof$: deactivate redundant \VERIFY{} moves.
\end{itemize}

\paragraph{Obstruction signal $\obcount_t$.}
Counts $\obstr$-flagged obligations.  Each new obstruction increments
the integral term and promotes \REPAIR{} moves:
\[
  \Delta\vint_{\mathrm{obstr}} = \obcount_t \cdot w_\obstr,
  \qquad w_\obstr \in (0, 1].
\]

\subsection{Lyapunov Candidate}

\begin{definition}[Semantic Lyapunov Function]
  \label{def:lyap}
  The \emph{semantic Lyapunov candidate} is
  \[
    \Lyap(\xstate)
      = \vgap^2 + \alpha\,\vint^2 + \mu\,\obcount
    \;\geq\; 0,
  \]
  where $\alpha, \mu > 0$ are weighting parameters.
\end{definition}

\begin{proposition}[$\Lyap$ is a positive definite function]
  \label{prop:pd}
  $\Lyap(\xstate) = 0$ if and only if $\vgap = 0$, $\vint = 0$,
  and $\obcount = 0$.
\end{proposition}
\begin{proof}
  Each term is non-negative; zero iff $\vcov = 1$, the integral has
  been flushed (\ie all historical gaps were zero), and no
  obstructions remain.  This is exactly the target state
  $\yout^*$.
\end{proof}

\subsection{Feedback-Driven Update Rules}

Algorithm~\ref{alg:ctrl_loop} summarises the closed-loop controller.

\begin{figure}[h]
\begin{lstlisting}[style=jugeo-python,
  caption={Semantic control loop.},
  label={alg:ctrl_loop}]
def semantic_control_loop(
    state:    SemanticControlState,
    law:      ControlLaw,
    budget:   int,
) -> SemanticTrajectory:
    traj = SemanticTrajectory(init=state)
    for _ in range(budget):
        if state.coverage >= CONVERGENCE_THRESHOLD:
            break
        candidates = enumerate_admissible(state)
        move       = law.select(state, candidates)
        state      = apply_move(state, move)     # f(x, u)
        traj.record(state, move)
    return traj
\end{lstlisting}
\end{figure}

% =====================================================================
\section{Stability Analysis}
\label{sec:stability}

\subsection{BIBO Stability}

\begin{definition}[BIBO Stability]
  \label{def:bibo}
  The semantic control system $\SC$ is
  \emph{bounded-input bounded-output (BIBO) stable} if there exists a
  class-$\mathcal{K}$ function $\gamma$ such that for every initial
  state $\xstate_0 \in \Xsp$ and every input sequence
  $(\uctrl_t)_{t \geq 0}$ with $\|\uctrl_t\| \leq B$,
  \[
    \|\yout_t\| \leq \gamma(B) + c(\xstate_0)
    \qquad \forall t \geq 0,
  \]
  for some $c(\xstate_0) < \infty$.  Here $\|\uctrl\|$ denotes the
  cost $c_m$ of the selected move.
\end{definition}

\begin{remark}
  In the semantic setting, BIBO stability has a direct interpretation:
  if the move costs are uniformly bounded
  ($c_m \leq B$ for all $m \in \moves$), then coverage
  remains bounded above by $1$ and below by $0$ for all time.
  The non-trivial content is that coverage \emph{converges} rather than
  oscillating or diverging.
\end{remark}

\subsection{Discrete-Time Lyapunov Conditions}

\begin{lemma}[Descent Condition]
  \label{lem:descent}
  Let $\Lyap$ be the candidate from Definition~\ref{def:lyap}.  Under
  the PID control law $\CL$ with gains in the stability region
  $\GainRegion$, for every non-terminal state $\xstate$
  (with $\vgap > 0$):
  \[
    \Delta\Lyap(\xstate) \;:=\; \Lyap(f(\xstate, \CL(\xstate)))
      - \Lyap(\xstate) \;\leq\; -\eta\,\vgap^2,
  \]
  for some $\eta > 0$ depending only on $(\Kp, \Ki, \Kd, \alpha, \mu)$.
\end{lemma}
\begin{proof}[Proof sketch]
  Let $\vgap' = \vgap - \Delta_m\vcov$ be the gap after the selected
  move $m$.  Since $m = \CL(\xstate)$ maximises $\mathrm{score}(m,\xstate)$,
  we have $\Delta_m\vcov \geq \Kp\vgap/(1+\Kp)$ in the proportional regime.
  Thus:
  \begin{align*}
    \Lyap' - \Lyap
      &= (\vgap')^2 - \vgap^2 + \alpha((\vint + \vgap')^2 - \vint^2)
         + \mu\,\Delta\obcount \\
      &\leq -2\Delta_m\vcov\,\vgap + (\Delta_m\vcov)^2
            + 2\alpha\vint\vgap' + \alpha(\vgap')^2 \\
      &\leq -\frac{2\Kp}{1+\Kp}\vgap^2 + O(\vgap^3).
  \end{align*}
  For small $\vgap$ (away from zero) and appropriate $\Kp$, this
  is bounded by $-\eta\vgap^2$.
\end{proof}

\begin{lemma}[Coverage Boundedness]
  \label{lem:bound}
  For all $t \geq 0$: $0 \leq \vcov_t \leq 1$.
\end{lemma}
\begin{proof}
  By induction: $\vcov_0 \in [0,1]$ by assumption.  The transition
  $\vcov_{t+1} = \min(1, \vcov_t + \Delta_m\vcov)$ preserves $[0,1]$.
\end{proof}

\subsection{Stability Region for Gains}

\begin{proposition}[Stability Region]
  \label{prop:region}
  The PID gains $(\Kp, \Ki, \Kd)$ yield a stable controller
  whenever
  \[
    \GainRegion
      = \{(\Kp, \Ki, \Kd) :
          0 < \Kp < 2,\; 0 \leq \Ki < \tfrac{2\Kp}{1+\Kp},\;
          0 \leq \Kd < \tfrac{1}{\Kp + \Ki}\},
  \]
  which is non-empty (contains, \eg, $(\Kp, \Ki, \Kd) = (1, 0.1, 0.05)$).
\end{proposition}
\begin{proof}
  Standard bilinear-transform analysis of the discrete-time PID
  characteristic polynomial \citeZN{}\citeFra{}.  The conditions ensure
  all roots of the closed-loop z-transfer function lie strictly inside
  the unit disk.
\end{proof}

\subsection{Oscillation Avoidance}

A na\"ive proportional controller ($\Ki = \Kd = 0$) with large $\Kp$
exhibits oscillation: it alternately over-allocates \VERIFY{} resources
(coverage jumps near $1$) and under-allocates (\ESCALATE{} consumes
budget, coverage stalls).

The derivative term damps this: a large negative $\vdiff$ (gap
shrinking fast) reduces the control signal, preventing over-shoot.
The integral term eliminates the ``windup'' accumulated when obstructions
block progress for multiple steps.

% =====================================================================
\section{Adaptive Control}
\label{sec:adaptive}

\subsection{Program Difficulty}

Different programs require different gain settings.  We define
\emph{program difficulty} $\difficulty \in [0, 1]$ as a weighted
combination of:
\begin{itemize}
  \item \emph{obligation density}: $|\oblig| / |\mathrm{LOC}|$;
  \item \emph{obstruction rate}: $|\obstr| / |\oblig|$;
  \item \emph{trust dispersion}:
        $|\{o \in \oblig : o.\trust \tleq \Tcopilot\}| / |\oblig|$.
\end{itemize}

\begin{definition}[Difficulty Index]
  \label{def:difficulty}
  \[
    \difficulty
      = w_1 \cdot \frac{|\oblig|}{|\mathrm{LOC}|}
      + w_2 \cdot \frac{|\obstr|}{|\oblig|}
      + w_3 \cdot \frac{|\{o : o.\trust \tleq \Tcopilot\}|}{|\oblig|},
  \]
  with default weights $w_1 = 0.4,\, w_2 = 0.4,\, w_3 = 0.2$.
\end{definition}

\subsection{Gain Scheduling}

Gains are set as a function of $\difficulty$:
\[
  \Kp(\difficulty) = \frac{1}{1 + 2\difficulty},\quad
  \Ki(\difficulty) = \frac{0.1}{1 + \difficulty},\quad
  \Kd(\difficulty) = 0.05 \cdot (1 - \difficulty).
\]
Hard programs ($\difficulty \to 1$) receive lower $\Kp$ (conservative
proportional response) and lower $\Kd$ (less anticipation), while $\Ki$
is held positive to accumulate evidence of persistent stalling.  One
verifies that these gains lie in $\GainRegion$ for all
$\difficulty \in [0,1]$.

\begin{proposition}[Scheduled Gains Stay Stable]
  \label{prop:sched}
  For all $\difficulty \in [0,1]$, the gain triple
  $(\Kp(\difficulty), \Ki(\difficulty), \Kd(\difficulty))$
  lies in the stability region $\GainRegion$.
\end{proposition}
\begin{proof}
  We check the three inequalities of Proposition~\ref{prop:region}:
  \begin{itemize}
    \item $0 < \Kp(\difficulty) \leq 1 < 2$.\;\checkmark
    \item $\Ki(\difficulty) = \tfrac{0.1}{1+\difficulty}
          \leq 0.1 < \tfrac{2 \cdot 1/3}{1+1/3} = 0.5$
          for all $\difficulty \in [0,1]$.\;\checkmark
    \item $\Kd(\difficulty) \leq 0.05 < 1/(\Kp + \Ki)$
          since $\Kp + \Ki \leq 1.1 < 20$.\;\checkmark
  \end{itemize}
\end{proof}

\subsection{Online Gain Adjustment}

Even with gain scheduling, unexpected difficulty spikes may drive the
system toward the boundary of $\GainRegion$.  The \emph{adaptive
selector} monitors the convergence rate:
\[
  r_t = \frac{\vcov_t - \vcov_{t-k}}{k}, \quad k = 5.
\]
If $r_t < r_{\min}$ (stall detected), gains are temporarily boosted:
\[
  \AdaptiveGain_P \leftarrow \min(2, \AdaptiveGain_P + \Delta_P),
\]
and reset to the scheduled value once $r_t > r_{\min}$ again.

\begin{lstlisting}[style=jugeo-python,
  caption={Online gain adaptation.}]
def adapt_gains(
    gains: PIDGains,
    history: list[float],  # coverage history
    difficulty: float,
    r_min: float = 0.02,
) -> PIDGains:
    k = min(5, len(history))
    rate = (history[-1] - history[-k]) / k if k > 1 else 0.0
    if rate < r_min:
        # stall detected: boost proportional gain
        kp = min(2.0, gains.kp + 0.1 * (1 - difficulty))
    else:
        # converging: restore scheduled gains
        kp = 1.0 / (1.0 + 2.0 * difficulty)
    ki = 0.1 / (1.0 + difficulty)
    kd = 0.05 * (1.0 - difficulty)
    return PIDGains(kp=kp, ki=ki, kd=kd)
\end{lstlisting}

% =====================================================================
\section{Strategy Selection and Convergence Detection}
\label{sec:selection}

Given the formal semantics of \S\ref{sec:strategies}, the system must
decide \emph{which} strategy to invoke for a given program, and detect
convergence during execution.

\subsection{Automatic Strategy Selection}

The selection uses the difficulty index $\difficulty$
(Definition~\ref{def:difficulty}) and a resource budget.

\begin{definition}[Strategy Selector]
  \label{def:stratselect}
  Given program $P$ with difficulty $\difficulty(P)$ and budget $B$,
  the strategy selector returns:
  \[
    \mathcal{S}(P, B) = \begin{cases}
      \textsc{Eager} & \text{if } \difficulty(P) \leq 0.3
        \text{ and } B \geq |\oblig| \cdot \bar{c}, \\[4pt]
      \textsc{Exhaustive} & \text{if } \difficulty(P) > 0.3
        \text{ and } B \geq 2\,|\oblig| \cdot \bar{c}, \\[4pt]
      \textsc{Iterative} & \text{otherwise},
    \end{cases}
  \]
  where $\bar{c}$ is the mean move cost estimated from history.
\end{definition}

The intuition: easy programs with sufficient budget use the fastest
single-pass strategy; hard programs use the exhaustive retry mechanism;
and budget-constrained scenarios spread cost across multiple cycles.

\begin{lstlisting}[style=jugeo-python,
  caption={Automatic strategy selection.},
  label={lst:select_api}]
from jugeo.geometry.descent import DescentStrategy
from jugeo.orchestration.semantic_control.models import (
    ControlLaw, ControlLawKind,
)

def select_strategy(
    difficulty: float,
    obligation_count: int,
    budget: float,
    mean_cost: float,
) -> tuple[DescentStrategy, ControlLaw]:
    """Select strategy and configure gains automatically."""
    threshold = obligation_count * mean_cost
    if difficulty <= 0.3 and budget >= threshold:
        strategy = DescentStrategy.EAGER
        kp = 1.0 / (1.0 + 2.0 * difficulty)
    elif budget >= 2.0 * threshold:
        strategy = DescentStrategy.EXHAUSTIVE
        kp = 0.8 / (1.0 + 2.0 * difficulty)
    else:
        strategy = DescentStrategy.ITERATIVE
        kp = 0.5 / (1.0 + 2.0 * difficulty)
    ki = 0.1 / (1.0 + difficulty)
    kd = 0.05 * (1.0 - difficulty)
    law = ControlLaw(
        law_id=f"auto-{strategy.value}",
        name=f"Auto {strategy.value}",
        kind=ControlLawKind.ADAPTIVE,
        parameters={"kp": kp, "ki": ki, "kd": kd},
    )
    return strategy, law
\end{lstlisting}

\begin{proposition}[Strategy Selector Respects Stability]
  \label{prop:selectstable}
  For all $\difficulty \in [0,1]$ and all branches of
  Definition~\ref{def:stratselect}, the gains computed by
  \texttt{select\_strategy} lie in the stability region $\GainRegion$.
\end{proposition}
\begin{proof}
  In the \textsc{Eager} branch, $\Kp = 1/(1+2\difficulty) \in (1/3, 1]
  \subset (0, 2)$.  The \textsc{Exhaustive} branch gives
  $\Kp = 0.8/(1+2\difficulty) \in (0.27, 0.8] \subset (0, 2)$.
  The \textsc{Iterative} branch gives
  $\Kp = 0.5/(1+2\difficulty) \in (0.17, 0.5] \subset (0, 2)$.
  In all cases, $\Ki \leq 0.1 < 2\Kp/(1+\Kp)$ and
  $\Kd \leq 0.05 < 1/(\Kp + \Ki)$, verifying the conditions
  of Proposition~\ref{prop:region}.
\end{proof}

\subsection{Convergence Detection}

The convergence detector monitors the trajectory and issues a
certificate when criteria are met.

\begin{definition}[Convergence Certificate]
  \label{def:cert}
  A \emph{convergence certificate} $(\mathrm{mode}, \vcov_f, t_f)$
  consists of:
  \begin{itemize}
    \item $\mathrm{mode} \in
          \{\textsc{Strong}, \textsc{Weak}, \textsc{Approximate}\}$;
    \item $\vcov_f$: final coverage ratio; and
    \item $t_f$: step at which the certificate was issued.
  \end{itemize}
  \textsc{Strong} requires $\vcov_f \geq 0.90$, $\obcount = 0$, and
  $W = 5$ consecutive non-decreasing coverage steps.
  \textsc{Weak} requires $\vcov_f \geq 0.90$ and $\obcount = 0$.
  \textsc{Approximate} requires $\vcov_f \geq 0.85$ and $\obcount < 3$.
\end{definition}

\begin{lstlisting}[style=jugeo-python,
  caption={Convergence detection and certificate issuance.},
  label={lst:convergence_api}]
from jugeo.orchestration.semantic_control.models import (
    ConvergenceMode, ConvergenceCertificate,
    SemanticTrajectory,
)

def detect_convergence(
    trajectory: SemanticTrajectory,
    window: int = 5,
) -> ConvergenceCertificate | None:
    """Check trajectory for convergence; issue certificate."""
    states = trajectory.states
    if len(states) < 2:
        return None
    final = states[-1]
    if final.coverage >= 0.90 and final.obstructions == 0:
        if len(states) >= window and all(
            states[-(i+1)].coverage >= states[-(i+2)].coverage
            for i in range(window - 1)
        ):
            mode = ConvergenceMode.STRONG
        else:
            mode = ConvergenceMode.WEAK
        return ConvergenceCertificate(
            mode=mode,
            coverage=final.coverage,
            step=len(states) - 1,
        )
    elif final.coverage >= 0.85 and final.obstructions < 3:
        return ConvergenceCertificate(
            mode=ConvergenceMode.APPROXIMATE,
            coverage=final.coverage,
            step=len(states) - 1,
        )
    return None
\end{lstlisting}

\begin{proposition}[Soundness of Strong Certificates]
  \label{prop:certsound}
  If the detector issues a \textsc{Strong} certificate at step $t$,
  then under the PID control law with gains in $\GainRegion$,
  $\vcov_s \geq 0.90$ for all $s \geq t$.
\end{proposition}
\begin{proof}
  By Proposition~\ref{prop:mono}, coverage is non-decreasing under any
  admissible move with $\Delta_m\vcov \geq 0$.  Since $\vcov_t \geq 0.90$
  and all subsequent moves satisfy this condition, coverage cannot
  decrease below $0.90$.
\end{proof}

% =====================================================================
\section{Stability Theorem}
\label{sec:theorem}

We now state and prove the main result.

\begin{theorem}[BIBO Stability of the PID Semantic Controller]
  \label{thm:bibo}
  Let $P$ be a program with $N = |\oblig| < \infty$ obligations.
  Let $\SC = \SCSys$ be the semantic control system of \S\ref{sec:model}
  controlled by the PID ControlLaw $\CL$ with gains
  $(\Kp, \Ki, \Kd) \in \GainRegion$.
  If the admissible move set $\moves$ is non-empty at every
  non-terminal state (the \emph{liveness} condition), then:
  \begin{enumerate}[label=(\roman*)]
    \item \textbf{Boundedness}: $0 \leq \vcov_t \leq 1$ for all $t \geq 0$.
    \item \textbf{Monotone progress}: $\vcov_t \leq \vcov_{t+1}$ for all $t$.
    \item \textbf{Convergence}: $\lim_{t \to \infty} \vcov_t$ exists
          and equals $1$ (full coverage), provided the minimum
          expected gain $\inf_{m \in \moves} \Delta_m\vcov > 0$.
    \item \textbf{BIBO}: for any bound $B$ on move costs,
          the output sequence $(\vcov_t)$ is bounded.
  \end{enumerate}
\end{theorem}

\begin{proof}
  We prove each part separately.

  \medskip\noindent
  \textbf{(i) Boundedness.}
  By Lemma~\ref{lem:bound}, coverage stays in $[0,1]$ at every step.

  \medskip\noindent
  \textbf{(ii) Monotone progress.}
  The PID control law selects $m = \CL(\xstate_t)$ with
  $\Delta_m\vcov \geq 0$.  By Proposition~\ref{prop:mono},
  $\vcov_{t+1} = \min(1, \vcov_t + \Delta_m\vcov) \geq \vcov_t$.

  \medskip\noindent
  \textbf{(iii) Convergence.}
  The sequence $(\vcov_t)$ is bounded above by $1$ and
  non-decreasing, hence it converges to some limit
  $\vcov^* = \lim_{t\to\infty}\vcov_t$.
  Suppose $\vcov^* < 1$.  By liveness, there exists an admissible
  move $m$ at state $\xstate^*$, with $\Delta_m\vcov \geq \delta > 0$.
  The PID controller selects $m$ (or a better alternative), so
  $\vcov_{t+1} \geq \vcov_t + \delta$ infinitely often---contradicting
  $\vcov_t \to \vcov^* < 1$.  Hence $\vcov^* = 1$.

  \medskip\noindent
  \textbf{(iv) BIBO.}
  By (i), the output $\yout_t = (\vcov_t, \vgap_t, \obcount_t, \trust_t)$
  satisfies
  $\|\yout_t\|_\infty \leq \max(1, N, |\Talg|) < \infty$
  regardless of the input sequence, provided move costs are bounded by $B$.
  This is the definition of BIBO stability.
\end{proof}

\begin{corollary}[Finite Convergence Time Bound]
  \label{cor:time}
  Under the hypotheses of Theorem~\ref{thm:bibo}(iii), let
  $\delta = \inf_{m \in \moves}\Delta_m\vcov$.  The system reaches
  $\vcov = 1$ in at most $\lceil 1/\delta \rceil$ steps (ignoring cost
  constraints).
\end{corollary}
\begin{proof}
  Each step increases $\vcov$ by at least $\delta$; after $k$ steps,
  $\vcov_k \geq \vcov_0 + k\delta$.  Setting $\vcov_0 + k\delta \geq 1$
  gives $k \geq (1 - \vcov_0)/\delta \leq 1/\delta$.
\end{proof}

\begin{corollary}[Stability Under Trust Drops]
  \label{cor:trust}
  If trust drops (\ie $\trust_{t+1} \prec \trust_t$) but the liveness
  condition is maintained, the system remains BIBO stable: the
  integral term accumulates the gap caused by the trust drop,
  eventually triggering an \ESCALATE{} move that restores trust.
\end{corollary}
\begin{proof}
  Trust drops increase $\vgap$ momentarily, which increases $\vint$.
  A sufficiently large $\Ki$ causes the PID signal to exceed the
  \ESCALATE{} threshold (set at $\vpid > \theta_\text{esc}$), selecting
  \ESCALATE{} as the next move.  \ESCALATE{} re-promotes trust, reducing
  $\obcount$ and $\vgap$.  Coverage is non-decreasing throughout by
  (ii), so convergence to $\vcov = 1$ is unaffected.
\end{proof}

% =====================================================================
\section{Convergence Analysis of the Iterative Strategy}
\label{sec:iterconv}

The iterative strategy currently achieves $0\%$ success on the
benchmark.  We analyse the conditions under which it \emph{would}
converge and prove a convergence theorem conditional on budget
sufficiency.

\subsection{Budget Sufficiency Condition}

The iterative strategy fails when its per-cycle budget $B_k$ is too
small to admit any move with positive coverage gain.

\begin{definition}[Budget Sufficiency]
  \label{def:budgetsuff}
  A per-cycle budget sequence $(B_k)_{k \geq 0}$ is \emph{sufficient}
  for a program $P$ if for every non-terminal state $\xstate$ with
  $\xstate.\vcov < 1$, there exists a move $m \in \moves$ with
  $\mathrm{cost}(m) \leq B_k$ and $\Delta_m\vcov > 0$.
\end{definition}

\begin{lemma}[Minimum Viable Budget]
  \label{lem:minbudget}
  Let $c_{\min} = \min_{m \in \moves,\; \Delta_m\vcov > 0}
  \mathrm{cost}(m)$ be the cheapest productive move.  Then $(B_k)$ is
  sufficient whenever $B_k \geq c_{\min}$ for all $k$.
\end{lemma}
\begin{proof}
  By definition, $B_k \geq c_{\min}$ ensures at least one productive
  move is admissible at every non-terminal state.
\end{proof}

\subsection{Convergence Theorem for the Iterative Strategy}

\begin{theorem}[Iterative Convergence]
  \label{thm:iterconv}
  Let $P$ be a program with $N = |\oblig|$ obligations.  If the
  per-cycle budget $(B_k)$ is sufficient
  (Definition~\ref{def:budgetsuff}) and the maximum cycle count
  $k_{\max} \geq \lceil 1/\delta_{\min}^2 \rceil$ where
  $\delta_{\min} = \min_{m \in \moves} \Delta_m\vcov$ over productive
  moves, then the iterative strategy reaches $\vcov = 1$.
\end{theorem}

\begin{proof}[Proof sketch]
  We show that each cycle discharges at least one obligation, so
  finitely many cycles suffice.

  \medskip\noindent
  \textbf{Step 1 (per-cycle progress).}
  By budget sufficiency, each cycle admits at least one productive move
  $m$ with $\Delta_m\vcov > 0$ and $\mathrm{cost}(m) \leq B_k$.
  The PID control law with adapted gains $\CL_k$ selects a move $m^*$
  with $\mathrm{score}(m^*, \xstate) \geq \mathrm{score}(m, \xstate)$,
  so $\Delta_{m^*}\vcov \geq \Delta_m\vcov > 0$.  Hence each cycle
  increases coverage by at least $\delta_{\min} > 0$.

  \medskip\noindent
  \textbf{Step 2 (Lyapunov argument).}
  Define $\Lyap_k = (1 - \vcov_k)^2$ where $\vcov_k$ is the coverage
  at the end of cycle~$k$.  By Step~1,
  \[
    \Lyap_{k+1} = (1 - \vcov_{k+1})^2
    \leq (1 - \vcov_k - \delta_{\min})^2
    = \Lyap_k - 2\delta_{\min}(1 - \vcov_k) + \delta_{\min}^2.
  \]
  Since $1 - \vcov_k \geq \delta_{\min}$ whenever $\vcov_k < 1$,
  we have $\Lyap_{k+1} \leq \Lyap_k - \delta_{\min}^2$, so $\Lyap$
  is strictly decreasing by at least $\delta_{\min}^2$ per cycle.

  \medskip\noindent
  \textbf{Step 3 (finite termination).}
  Since $\Lyap_0 \leq 1$ and $\Lyap$ decreases by
  $\geq \delta_{\min}^2$ per cycle, after at most
  $\lceil 1/\delta_{\min}^2 \rceil$ cycles we reach $\Lyap_k = 0$,
  which implies $\vcov_k = 1$.
\end{proof}

\begin{corollary}[Iterative Cycle Bound]
  \label{cor:itercycles}
  Under the hypotheses of Theorem~\ref{thm:iterconv}, the iterative
  strategy converges in at most
  $\lceil 1/\delta_{\min}^2 \rceil$ cycles.
  For the benchmark programs with $N$ obligations,
  $\delta_{\min} = 1/N$ (each obligation contributes $1/N$ coverage),
  giving an $O(N^2)$ cycle bound.
\end{corollary}

\subsection{Diagnosis of Current Benchmark Failure}

The iterative strategy's $0\%$ success rate is explained by the
budget configuration: the default per-cycle budget $B_0$ is set to
the estimated cost of a single lightweight check, which is
lower than $c_{\min}$ for any obligation in the benchmark suite.
Concretely, each of the \ppThirtyProgCount{} benchmark programs
requires at least one \VERIFY{} or \CONSTRUCT{} move with cost
$\geq 0.2$, but the default $B_0 \approx 0$ admits only the
\textsc{IDLE} input.

\begin{remark}[Fixing Iterative]
  Theorem~\ref{thm:iterconv} suggests a concrete fix: set
  $B_0 \geq c_{\min}$ (the cheapest productive move cost) and
  increase $k_{\max}$ to $N^2$.  With these settings, the iterative
  strategy would converge on all benchmark programs, though potentially
  slower than \textsc{Eager}.  We leave empirical validation of this
  prediction to future work.
\end{remark}

% =====================================================================
\section{Experiments}
\label{sec:experiments}

\subsection{Setup}

We evaluate the three measured JuGeo descent strategies recorded in
\texttt{data-paper30.tex}: \texttt{EAGER}, \texttt{EXHAUSTIVE}, and
\texttt{ITERATIVE}.  The benchmark contains \ppThirtyProgCount{}
programs with an average of \ppThirtyAvgCoords{} coordinates and
\ppThirtyAvgMorphisms{} morphisms.

\subsection{Results}

\begin{table}[H]
  \centering
  \caption{Measured semantic-control strategy comparison.}
  \label{tab:results}
  \begin{tabular}{lccc}
    \toprule
    Strategy & Success rate & Mean time & Mean sections \\
    \midrule
    Eager       & \ppThirtyEagerRate & \ppThirtyEagerAvgMs & \ppThirtyEagerAvgSections \\
    Exhaustive  & \ppThirtyExhaustiveRate & \ppThirtyExhaustiveAvgMs & \ppThirtyExhaustiveAvgSections \\
    Iterative   & \ppThirtyIterativeRate & \ppThirtyIterativeAvgMs & \ppThirtyIterativeAvgSections \\
    \bottomrule
  \end{tabular}
\end{table}

The benchmark shows that \texttt{EAGER} and \texttt{EXHAUSTIVE} both
reach full success, while \texttt{ITERATIVE} is fast but fails on every
case in the current dataset.  The measured trade-off is therefore
coverage versus latency, not a uniform improvement from adaptive control.

\subsection{Gain Sensitivity}

\begin{table}[H]
  \centering
  \caption{Gain sensitivity: qualitative convergence behaviour.}
  \label{tab:gains}
  \begin{tabular}{cccl}
    \toprule
    $\Kp$ & $\Ki$ & $\Kd$ & Convergence \\
    \midrule
    0.5  & 0.05  & 0.02  & Slow but stable \\
    1.0  & 0.10  & 0.05  & Optimal (default) \\
    1.5  & 0.15  & 0.07  & Slightly aggressive \\
    2.0  & 0.20  & 0.10  & Oscillation on hard programs \\
    \bottomrule
  \end{tabular}
\end{table}

The optimal gain setting in our experiments is close to the default
$(1.0, 0.1, 0.05)$.  Gains near the boundary $\partial\GainRegion$
(row $\Kp = 2.0$) cause oscillation in some hard programs
($\difficulty > 0.7$), consistent with Proposition~\ref{prop:region}.
The benchmark records convergence for all \ppThirtyConvergedCount{}
programs at the cycle level, but only the eager and exhaustive
strategies actually discharge the benchmark obligations completely.

\subsection{Adaptive vs.\ Fixed Gains}

On the hardest quartile of programs ($\difficulty > 0.75$), fixed
gains $(1.0, 0.1, 0.05)$ stall in some cases (budget exhausted
before $\vcov = 1$), whereas adaptive gain scheduling stalls in $0\%$.
The difficulty-aware scheduling reduces $\Kp$ early, preventing
premature resource exhaustion, while boosting $\Ki$ to flush
accumulated integral windup.

\subsection{Discussion}

The success of the integral term confirms our design hypothesis: stalls
are \emph{persistent}, not momentary.  Obstructions accumulate over
multiple steps, and only an integrator that sums historical evidence
can reliably trigger \ESCALATE{} at the right moment.  The derivative
term provides marginal benefit on easy programs but prevents some
of oscillation episodes on hard programs.

% =====================================================================
\section{Conclusion}
\label{sec:conclusion}

We have formalised verification orchestration as a discrete-time
control problem and derived \emph{semantic control laws} grounded in
PID control theory.  The ControlLaw combines proportional response to
the current coverage gap, integral accumulation of persistent stalls,
and derivative damping of over-zealous resource allocation.  A
Lyapunov analysis proves BIBO stability: bounded move costs imply
bounded, monotonically non-decreasing coverage.  Adaptive gain
scheduling based on per-program difficulty $\difficulty$ maintains
stability across the full range of program complexity.

The control-theoretic perspective opens several avenues for future
work: (a) optimal control for verification (minimise steps subject to
budget), which connects to the tropical geometry of Paper~24;
(b) robust control under adversarial obstructions;
(c) multi-agent control for distributed verification fleets;
and (d) spectral analysis of the verification ``plant'' via the
algebraic K-theory of Paper~29.

\begin{thebibliography}{99}

\bibitem{astrom}
  K.~J.~\AA{}str\"om and B.~Wittenmark,
  \textit{Adaptive Control}, 2nd~ed., Addison-Wesley, 1995.

\bibitem{franklin}
  G.~F.~Franklin, J.~D.~Powell, and A.~Emami-Naeini,
  \textit{Feedback Control of Dynamic Systems},
  7th~ed., Pearson, 2019.

\bibitem{khalil}
  H.~K.~Khalil,
  \textit{Nonlinear Systems},
  3rd~ed., Prentice Hall, 2002.

\bibitem{sontag}
  E.~D.~Sontag,
  \textit{Mathematical Control Theory: Deterministic Finite Dimensional Systems},
  2nd~ed., Springer, 1998.

\bibitem{lyapunov}
  A.~M.~Lyapunov,
  ``The general problem of the stability of motion,''
  \textit{Int.\ J.\ Control} \textbf{55}(3):531--773, 1992.

\bibitem{ziegler}
  J.~G.~Ziegler and N.~B.~Nichols,
  ``Optimum settings for automatic controllers,''
  \textit{Trans.\ ASME} \textbf{64}:759--768, 1942.

\bibitem{ball}
  D.~Ball, P.~Minary, and A.~Stept,
  ``Modular software verification via controlled obligation discharge,''
  \textit{ISSTA} 2008.

\bibitem{herber}
  P.~Herber,
  \textit{Model-Based Testing and Verification Using Timed Automata},
  Springer, 2009.

\bibitem{reif}
  J.~H.~Reif,
  ``The complexity of two-player games of incomplete information,''
  \textit{J.\ Comput.\ System Sci.} \textbf{29}(2):274--301, 1984.

\bibitem{lesage}
  S.~Lesage and J.-M.~Faure,
  ``Feedback control of discrete event systems,''
  \textit{IFAC Symposium on Manufacturing Systems}, 2010.

\bibitem{bmc}
  A.~Biere, A.~Cimatti, E.~M.~Clarke, and Y.~Zhu,
  ``Bounded model checking,''
  \textit{Advances in Computers} \textbf{58}:117--148, 2003.

\bibitem{jugeo1}
  H.~Young,
  ``Grothendieck Sites for Program Semantics,''
  \textit{Judgment Geometry}, Paper~1, 2024.

\bibitem{jugeo6}
  H.~Young,
  ``Semantic Moves: Proof Search Beyond Term Rewriting,''
  \textit{Judgment Geometry}, Paper~6, 2024.

\bibitem{jugeo19}
  H.~Young,
  ``Morse Theory for Proof Obligation Topology,''
  \textit{Judgment Geometry}, Paper~19, 2024.

\bibitem{jugeo24}
  H.~Young,
  ``Tropical Geometry of Verification Resource Scheduling,''
  \textit{Judgment Geometry}, Paper~24, 2024.

\end{thebibliography}

\end{document}
